\documentclass[11pt]{article}

\usepackage{amsmath}
\usepackage{amssymb}
\usepackage{booktabs}
\usepackage[margin=1in]{geometry}
\usepackage{hyperref}
\hypersetup{colorlinks=true,linkcolor=blue,citecolor=blue,urlcolor=blue}

\newcommand{\vhat}{\hat{\mathbf{v}}}
\newcommand{\ig}{\hat{\boldsymbol{\imath}}_{\gamma}}
\newcommand{\ip}{\hat{\boldsymbol{\imath}}_{\psi}}
\newcommand{\nhat}{\hat{\mathbf{n}}}
\newcommand{\ehat}{\hat{\mathbf{e}}}
\newcommand{\uhat}{\hat{\mathbf{u}}}
\newcommand{\zhat}{\hat{\mathbf{z}}}
\newcommand{\bomega}{\boldsymbol{\omega}}
\newcommand{\bOmega}{\boldsymbol{\Omega}}
\newcommand{\code}[1]{\texttt{#1}}
\newcommand{\dd}{\mathrm{d}}

\title{Boost, Glide and Descent Dynamics\\
       \large The mathematics of the \code{missiles/} trajectory library}
\author{Coorbital, Inc.}
\date{7 August 2026}

\begin{document}
\maketitle

\begin{abstract}
This note derives the three-degree-of-freedom point-mass equations of motion
used by the \code{missiles/} MATLAB library for powered boost, unpowered
atmospheric glide and terminal descent over a spherical, optionally rotating
Earth; the four injected environment models; the phase machinery that chains
them; and the seven closed-form references the library is validated against.
It is written to be implementable from: every equation here is the equation the
code evaluates, checked term by term against the source, and every closed form
is derived rather than quoted. Where a result has a numerical consequence, the
measured value is given. Symbolic checks were run in \code{sympy}~1.14 and
numerical checks in MATLAB R2025b against the committed library;
Section~\ref{sec:verify} records what was verified and how.
\end{abstract}

\tableofcontents

%=====================================================================
\section{Scope, and how to read this note}

\subsection{What is here}

The library flies a point mass over a sphere. Its state is six numbers ---
position and velocity in geocentric spherical coordinates --- with a seventh,
mass, appended whenever a chain contains a powered phase. Everything in
Sections~\ref{sec:frames}--\ref{sec:phases} is a derivation of what
\code{+coorbital/+eom/glide3DOF.m} and \code{+coorbital/+eom/boost3DOF.m}
actually compute, in the order they compute it.
Section~\ref{sec:analytic} derives the analytic solutions the library is
measured against and reports the measured agreement.
Section~\ref{sec:targeting} covers the two-solve targeting problem.

\subsection{What is verified, and how}
\label{sec:verify}

Three kinds of check stand behind the statements below.

\begin{description}
\item[Symbolic.] The velocity triad's orthonormality and handedness, the three
gravity projections, the three thrust projections, the two transport
(curvature) terms, and the Coriolis and centrifugal projections were all
re-derived in \code{sympy}~1.14 from explicit vector cross products and
compared against the expressions in the source. All matched identically.
\item[Numerical, against the code.] Each projection was independently
reconstructed in MATLAB from explicit north--east--up basis vectors and
compared against the derivative the library returns. Worst relative residual
over all channels: $9.5\times10^{-16}$ for the gravity and aerodynamic terms,
$1.6\times10^{-16}$ for thrust, $1.1\times10^{-13}$ for the rotating-Earth
terms. Critically, this was done with a \emph{synthetic} gravity model
returning $g_{\phi}=0.0123~\mathrm{m/s^2}$, because the shipped
\code{sphereGrav} returns $g_{\phi}\equiv0$ and therefore cannot exercise the
latitudinal terms at all.
\item[Numerical, against closed forms.] Each analytic reference in
Section~\ref{sec:analytic} was re-measured on the committed code. The numbers
quoted are from that run, not carried over from a report.
\end{description}

\subsection{What this note is not}

It is not a validation report --- \code{README.md} is the authority on
delivered behaviour, and \code{DESIGN.md} records where the
pre-implementation spec and the code disagree. It does not cover
aerothermodynamics: no heating model exists in the library. It does not cover
6-DOF attitude dynamics; angle of attack and bank angle are \emph{commanded}
here, not the output of a rotational state. And every vehicle and booster
number that appears below is a marked open-literature \textsc{placeholder};
they are of the right order and are not any specific system.

%=====================================================================
\section{Frames, coordinates and conventions}
\label{sec:frames}

\subsection{The state}

Position is geocentric spherical: radius from the planet centre, longitude,
geocentric latitude. Velocity is planet-relative and given by its magnitude and
two angles. The integrated state is
\begin{equation}
  \mathbf{x} \;=\; \begin{bmatrix} r & \lambda & \phi & V & \gamma & \psi
  \end{bmatrix}^{\!\top},
  \label{eq:state}
\end{equation}
with the total mass currently carried appended as a seventh component,
$\mathbf{x}_7 = m$, on any chain that contains a powered phase
(Section~\ref{sec:phases}). Four conventions are fixed once and never varied:

\begin{itemize}
\item $\gamma$ is \textbf{positive up}. A climbing vehicle has
$\gamma>0$; a descending one has $\gamma<0$. Impact flight-path angles are
therefore negative throughout the library.
\item $\psi$ is measured \textbf{clockwise from north}: $0$ is due north,
$\pi/2$ due east. This is deliberately the same convention as a navigational
bearing, so \code{coorbital.util.greatCircleBearing} drops straight into
$\mathbf{x}_6$ with no sign flip and no offset
(Section~\ref{sec:targeting}).
\item $V$ is \textbf{planet-relative}, not inertial. With the Earth rate
switched on, $V$ is the speed seen by an observer rotating with the planet, and
the Coriolis and centrifugal terms of Section~\ref{sec:rotating} are what make
that consistent.
\item $\phi$ is \textbf{geocentric} latitude, and $h = r - r_E$ is altitude
above a sphere of radius $r_E$. There is no ellipsoid anywhere in the library,
so geocentric and geodetic latitude are the same thing here; they will not be
when an oblate model arrives.
\end{itemize}

\subsection{The local horizon basis}

At the sub-vehicle point define the usual north--east--up triple
$(\nhat,\ehat,\uhat)$: $\uhat$ along the outward radius, $\nhat$ tangent to the
meridian toward the north pole, $\ehat$ completing the set to the east. Written
in that \emph{listed} order the triple is left-handed; the right-handed
ordering is $(\ehat,\nhat,\uhat)$, and the cross products used throughout are
\begin{equation}
  \ehat\times\nhat = \uhat, \qquad
  \nhat\times\uhat = \ehat, \qquad
  \uhat\times\ehat = \nhat .
  \label{eq:handed}
\end{equation}
This is worth stating explicitly because getting it backwards flips the sign of
every Coriolis term at once, and the resulting trajectory still looks
plausible. The planet's spin axis resolves in this basis as
\begin{equation}
  \zhat \;=\; \cos\phi\,\nhat + \sin\phi\,\uhat ,
  \label{eq:zhat}
\end{equation}
which is the only place the latitude enters the rotation terms.

\subsection{The velocity triad}

The three angles $(V,\gamma,\psi)$ are the spherical components of the
planet-relative velocity in that basis. Writing $\mathbf{V} = V\vhat$,
\begin{align}
  \vhat &= \cos\gamma\cos\psi\;\nhat + \cos\gamma\sin\psi\;\ehat
           + \sin\gamma\;\uhat ,   \label{eq:vhat}\\
  \ig   &= -\sin\gamma\cos\psi\;\nhat - \sin\gamma\sin\psi\;\ehat
           + \cos\gamma\;\uhat ,   \label{eq:igam}\\
  \ip   &= -\sin\psi\;\nhat + \cos\psi\;\ehat .  \label{eq:ipsi}
\end{align}
$\vhat$ points along the velocity; $\ig$ is the direction in which a positive
increment of $\gamma$ moves the velocity, i.e.\ ``up and normal to the flight
path''; $\ip$ is the direction in which a positive increment of $\psi$ moves
it, i.e.\ ``horizontally to the right of the track''.

\paragraph{The triad is orthonormal and right-handed.}
Direct computation from \eqref{eq:vhat}--\eqref{eq:ipsi}:
\begin{align}
  \vhat\cdot\vhat &= \cos^2\!\gamma(\cos^2\!\psi+\sin^2\!\psi) + \sin^2\!\gamma
                   = \cos^2\!\gamma + \sin^2\!\gamma = 1, \nonumber\\
  \ig\cdot\ig   &= \sin^2\!\gamma + \cos^2\!\gamma = 1, \qquad
  \ip\cdot\ip    = \sin^2\!\psi + \cos^2\!\psi = 1, \nonumber\\
  \vhat\cdot\ig &= -\sin\gamma\cos\gamma(\cos^2\!\psi+\sin^2\!\psi)
                   + \sin\gamma\cos\gamma = 0, \nonumber\\
  \vhat\cdot\ip &= -\cos\gamma\cos\psi\sin\psi + \cos\gamma\sin\psi\cos\psi = 0,
                   \nonumber\\
  \ig\cdot\ip   &= \phantom{-}\sin\gamma\cos\psi\sin\psi
                   - \sin\gamma\sin\psi\cos\psi = 0 .
  \label{eq:orthonormal}
\end{align}
Every cross term collapses on the same Pythagorean identity, which is the
structural reason the triad is orthonormal for \emph{all} $(\gamma,\psi)$ and
not merely for small angles. The handedness follows from \eqref{eq:handed}:
$\vhat\times\ig = \ip$, verified symbolically. Numerically the Gram matrix of
the triad differs from the identity by $1.1\times10^{-16}$ at a general state.

Two derivatives that will be needed immediately, and which are the whole reason
this particular triad is chosen:
\begin{equation}
  \frac{\partial\vhat}{\partial\gamma} = \ig ,
  \qquad
  \frac{\partial\vhat}{\partial\psi} = \cos\gamma\;\ip .
  \label{eq:dvhat}
\end{equation}
The $\cos\gamma$ in the second is the first appearance of the factor that will
haunt the heading equation: near vertical flight a large change of heading
barely moves the velocity vector, so a given side force produces a very large
$\dot\psi$.

\subsection{Table of symbols}

\begin{center}
\begin{tabular}{llll}
\toprule
Symbol & Meaning & Units & Code \\
\midrule
$r$        & geocentric radius                        & m         & \code{x(1)} \\
$\lambda$  & longitude                                & rad       & \code{x(2)} \\
$\phi$     & geocentric latitude                      & rad       & \code{x(3)} \\
$V$        & planet-relative speed                    & m/s       & \code{x(4)} \\
$\gamma$   & flight-path angle, positive up           & rad       & \code{x(5)} \\
$\psi$     & heading, clockwise from north            & rad       & \code{x(6)} \\
$m$        & total mass currently carried             & kg        & \code{x(7)} \\
$\alpha$   & angle of attack                          & rad       & \code{u(1)} \\
$\sigma$   & bank angle about the velocity vector     & rad       & \code{u(2)} \\
$h$        & altitude above the reference sphere      & m         & \code{h} \\
$g_r$      & inward gravitational attraction, \emph{positive magnitude} & m/s$^2$ & \code{gr} \\
$g_\phi$   & northward gravitational acceleration     & m/s$^2$   & \code{gLat} \\
$\rho$     & atmospheric density                      & kg/m$^3$  & \code{rho} \\
$P$        & ambient static pressure                  & Pa        & \code{P} \\
$a_s$      & speed of sound                           & m/s       & \code{aSnd} \\
$\bar q$   & dynamic pressure $\tfrac12\rho V^2$      & Pa        & \code{qbar} \\
$S$        & aerodynamic reference area               & m$^2$     & \code{veh.Sref} \\
$C_L,C_D$  & lift and drag coefficients               & --        & \code{CL,CD} \\
$L,D$      & lift and drag forces                     & N         & --- \\
$T$        & delivered thrust                         & N         & \code{T} \\
$\dot m_p$ & propellant mass flow, \emph{positive}    & kg/s      & \code{mdot} \\
$\mu$      & gravitational parameter                  & m$^3$/s$^2$ & \code{c.muE} \\
$r_E$      & reference sphere radius                  & m         & \code{c.rE} \\
$\omega$   & planet rotation rate                     & rad/s     & \code{env.omegaE} \\
$H$        & density scale height                     & m         & \code{c.Hscale} \\
$I_{sp}$   & vacuum specific impulse                  & s         & \code{bst.Isp} \\
$g_0$      & standard gravity, $9.80665$              & m/s$^2$   & \code{c.g0} \\
$\beta$    & ballistic coefficient $m/(SC_D)$         & kg/m$^2$  & --- \\
\bottomrule
\end{tabular}
\end{center}

%=====================================================================
\section{Kinematics}
\label{sec:kinematics}

The position vector is $\mathbf{r} = r\,\uhat(\lambda,\phi)$ and the velocity
is $\mathbf{V}=V\vhat$. Resolving the velocity on the horizon basis and
matching it to the time derivative of the position gives the three kinematic
equations directly.

The radial component is immediate:
\begin{equation}
  \dot r \;=\; \mathbf{V}\cdot\uhat \;=\; V\sin\gamma .
  \label{eq:rdot}
\end{equation}

The two horizontal components must be matched against the surface velocity that
a change in $(\lambda,\phi)$ produces. Moving north at fixed longitude sweeps
arc length $r\,\dd\phi$; moving east at fixed latitude sweeps arc length
$r\cos\phi\,\dd\lambda$, because the parallel of latitude $\phi$ is a circle of
radius $r\cos\phi$, not $r$. Hence
\begin{equation}
  r\dot\phi = \mathbf{V}\cdot\nhat = V\cos\gamma\cos\psi ,
  \qquad
  r\cos\phi\;\dot\lambda = \mathbf{V}\cdot\ehat = V\cos\gamma\sin\psi ,
\end{equation}
that is
\begin{equation}
  \dot\phi \;=\; \frac{V\cos\gamma\cos\psi}{r},
  \qquad
  \dot\lambda \;=\; \frac{V\cos\gamma\sin\psi}{r\cos\phi} .
  \label{eq:latlondot}
\end{equation}
These are exactly the three lines under \code{\%\% Kinematics:} in both EOM
files.

\paragraph{Why $\cos\phi$ sits in the denominator, and what breaks at the pole.}
The $\cos\phi$ is a metric factor, not a physical one. Longitude is an angle
measured about the polar axis, and near the pole a small eastward
\emph{distance} corresponds to a large change of longitude, because the
parallel it is measured on has shrunk to nothing. At the pole itself
$\cos\phi=0$ and $\dot\lambda$ is unbounded: the coordinate is degenerate, not
the physics. A vehicle flying over the pole has a perfectly well-behaved
velocity vector; it is only its \emph{longitude} that becomes meaningless and
then jumps by $\pi$.

An adaptive integrator meeting this does not fail cleanly. It sees a stiff
right-hand side, cuts its step, and either grinds to a halt or produces a
longitude history that is numerically arbitrary while every other state looks
fine. Both EOM files therefore guard it explicitly, raising
\code{coorbital:glide3DOF:polarSingularity} when $|\cos\phi|<10^{-8}$, rather
than letting a silent \code{NaN} propagate into the trajectory. The
corresponding structural fix --- a non-singular parameterisation, or a
Cartesian terminal phase --- is not in the library and is not needed for the
mid-latitude trajectories it flies.

\paragraph{A consistency reading.} Set $\gamma=0$ in \eqref{eq:rdot}: the
radius is stationary, as it must be for horizontal flight. Set $\psi=\pi/2$
(due east) at the equator in \eqref{eq:latlondot}: $\dot\phi=0$ and
$\dot\lambda>0$. The library asserts both of these directly
(\code{test\_glide3DOF.m}), and asserts the $\cos\phi$ placement by checking
$\dot\lambda\,r\cos\phi = V\cos\gamma\sin\psi$ to $10^{-10}$ relative at an
off-equator state --- which is the assertion that a $\cos\phi$ moved to the
numerator could not survive.

%=====================================================================
\section{Dynamics}
\label{sec:dynamics}

\subsection{The general form}

Differentiate $\mathbf{V}=V\vhat$ once:
\begin{equation}
  \mathbf{a} \;=\; \dot V\,\vhat + V\,\frac{\dd\vhat}{\dd t} .
  \label{eq:adecomp}
\end{equation}
The unit vector $\vhat$ changes for two distinct reasons: the angles
$(\gamma,\psi)$ change, and the basis $(\nhat,\ehat,\uhat)$ itself rotates as
the vehicle moves over the sphere. Writing $\bomega_f$ for the angular velocity
of the horizon basis \textbf{relative to the planet-fixed rotating frame} and
using \eqref{eq:dvhat},
\begin{equation}
  \frac{\dd\vhat}{\dd t}
  \;=\; \dot\gamma\,\ig + \dot\psi\cos\gamma\,\ip + \bomega_f\times\vhat .
  \label{eq:dvhatdt}
\end{equation}

\paragraph{Which frame $\bomega_f$ is measured in, and why it is not inertial.}
This is load-bearing and is fixed here once. $\mathbf{V} = V\vhat$ is the
\emph{planet-relative} velocity by the convention of
Section~\ref{sec:frames}, so \eqref{eq:adecomp} is a derivative taken in the
planet-fixed rotating frame, and $\bomega_f$ has to be the basis rate seen from
that same frame. It therefore contains the vehicle's motion over the sphere and
nothing else --- which is exactly what \eqref{eq:omegaf} below turns out to be,
built from $\dot\lambda$ and $\dot\phi$ with no $\omega$ in it. The planet's own
spin is deliberately absent: the inertial rate of the horizon basis is
$\bOmega + \bomega_f$, and the $\bOmega$ part enters the equations of motion
instead as the Coriolis and centrifugal accelerations of
Section~\ref{sec:rotating}. Taking $\bomega_f$ inertial \emph{and} adding those
pseudo-forces would count the planet's rotation twice; taking it inertial and
omitting them would be the other consistent choice, and is not the one the code
makes.

Projecting \eqref{eq:adecomp} on the triad and using orthonormality
\eqref{eq:orthonormal} --- $\vhat\cdot\dd\vhat/\dd t = 0$ for a unit vector ---
gives three scalar equations:
\begin{align}
  \dot V      &= \mathbf{a}\cdot\vhat ,                                        \label{eq:Vgen}\\
  \dot\gamma  &= \frac{\mathbf{a}\cdot\ig}{V} - (\bomega_f\times\vhat)\cdot\ig ,\label{eq:gamgen}\\
  \dot\psi    &= \frac{\mathbf{a}\cdot\ip}{V\cos\gamma}
                 - \frac{(\bomega_f\times\vhat)\cdot\ip}{\cos\gamma} .          \label{eq:psigen}
\end{align}

\paragraph{The transport terms.}
Seen from the rotating frame the horizon basis turns only because the vehicle
moves over it --- and that is true whether or not the planet spins, which is
what makes the block below independent of $\omega$. Increasing longitude rotates the basis about the polar axis at
$\dot\lambda$; increasing latitude tilts it about the \emph{westward}
horizontal, since $\dd\uhat/\dd t = \dot\phi\,\nhat$ and, by
\eqref{eq:handed}, $\ehat\times\uhat = -\nhat$. Hence
\begin{equation}
  \bomega_f \;=\; \dot\lambda\,\zhat - \dot\phi\,\ehat .
  \label{eq:omegaf}
\end{equation}
Substituting \eqref{eq:zhat}, \eqref{eq:latlondot} and \eqref{eq:vhat} and
reducing (verified symbolically) gives
\begin{equation}
  (\bomega_f\times\vhat)\cdot\vhat = 0,
  \qquad
  (\bomega_f\times\vhat)\cdot\ig = -\frac{V\cos\gamma}{r},
  \qquad
  (\bomega_f\times\vhat)\cdot\ip = -\frac{V\cos^2\!\gamma\sin\psi\tan\phi}{r} .
  \label{eq:transport}
\end{equation}
Putting \eqref{eq:transport} into \eqref{eq:gamgen}--\eqref{eq:psigen}:
\begin{align}
  \dot V     &= \mathbf{a}\cdot\vhat ,                                             \label{eq:Vgen2}\\
  \dot\gamma &= \frac{\mathbf{a}\cdot\ig}{V} + \frac{V}{r}\cos\gamma ,              \label{eq:gamgen2}\\
  \dot\psi   &= \frac{\mathbf{a}\cdot\ip}{V\cos\gamma}
                + \frac{V}{r}\cos\gamma\sin\psi\tan\phi .                           \label{eq:psigen2}
\end{align}
These three, with \eqref{eq:rdot} and \eqref{eq:latlondot}, are the whole
model. Everything that follows is a statement about what goes into
$\mathbf{a}$.

\subsection{Gravity}
\label{sec:gravity}

The gravity model interface returns two numbers, $(g_r,g_\phi)$, and the
convention on the first is the single most likely source of a sign error in the
whole library:

\begin{quote}
\textbf{$g_r$ is a positive magnitude of \emph{inward} attraction, not a signed
radial component.} For point-mass gravity $g_r = \mu/r^2 > 0$ always. The minus
sign that turns it into a downward acceleration lives in the equations of
motion, not in the model.
\end{quote}

The reason for the convention is interface stability rather than elegance: a
scalar ``local up'' value cannot represent $J_2$, whose acceleration has a
latitudinal component, so the interface commits to two components from the
start and \code{sphereGrav} returns $g_\phi\equiv 0$. Written as a vector on the
horizon basis,
\begin{equation}
  \mathbf{g} \;=\; g_\phi\,\nhat + 0\,\ehat - g_r\,\uhat .
  \label{eq:gvec}
\end{equation}
Its three projections follow from \eqref{eq:vhat}--\eqref{eq:ipsi} by
inspection, each being a single dot product:
\begin{align}
  \mathbf{g}\cdot\vhat &= -g_r\sin\gamma + g_\phi\cos\gamma\cos\psi ,
    \label{eq:gv}\\
  \mathbf{g}\cdot\ig   &= -g_r\cos\gamma - g_\phi\sin\gamma\cos\psi ,
    \label{eq:gg}\\
  \mathbf{g}\cdot\ip   &= -g_\phi\sin\psi .
    \label{eq:gp}
\end{align}
All three were re-derived symbolically and match the source term for term.

\paragraph{Physical reading, and the cost of each sign.}
Equation~\eqref{eq:gv} says gravity removes speed when climbing
($\gamma>0 \Rightarrow -g_r\sin\gamma<0$) and adds it when descending. Flip that
sign and a descending vehicle decelerates under gravity: on a ballistic entry
the impact speed collapses, and, worse, the vacuum energy check of
Section~\ref{sec:energy} fails immediately, which is precisely why that check
is the library's first line of defence.

Equation~\eqref{eq:gg} says gravity always bends the flight path \emph{down},
by $-g_r\cos\gamma/V$, and does so most strongly in level flight and not at all
in vertical flight --- a vertically-moving vehicle is already pointed along
gravity, so gravity can change its speed but not its direction. Flip this sign
and the vehicle lofts under its own weight; the trajectory is obviously wrong
within seconds.

Equation~\eqref{eq:gp} is the only term in which a purely northward gravity
component can steer the heading, and it is identically zero for a point mass.
It is carried anyway so that dropping in \code{j2Grav} requires no edit to the
equations of motion. Because $g_\phi\equiv 0$ under the shipped model, the
$g_\phi$ terms in \eqref{eq:gv}--\eqref{eq:gp} are \emph{never exercised} by
the shipped test suite. The verification for this note therefore injected a
synthetic gravity model with $g_\phi=0.0123~\mathrm{m/s^2}$ and confirmed all
three projections against an independent vector reconstruction to
$9.5\times10^{-16}$ relative. A reviewer should treat those three terms as
verified by this note and by nothing else in the repository.

\subsection{Aerodynamics}

Lift and drag are the usual
\begin{equation}
  L = \bar q\,S\,C_L , \qquad D = \bar q\,S\,C_D , \qquad
  \bar q = \tfrac12\rho V^2 .
  \label{eq:LD}
\end{equation}
Drag opposes the velocity, so it is purely along $-\vhat$. Lift is by
definition perpendicular to the velocity, which leaves it one degree of
freedom: the angle by which it is rolled about the velocity vector. That angle
is the bank angle $\sigma$, and it distributes the lift between the two normal
channels:
\begin{equation}
  \mathbf{a}_{\text{aero}} \;=\;
  -\frac{D}{m}\,\vhat \;+\; \frac{L}{m}\bigl(\cos\sigma\,\ig
  + \sin\sigma\,\ip\bigr) .
  \label{eq:aaero}
\end{equation}
Hence, with \eqref{eq:Vgen2}--\eqref{eq:psigen2},
\begin{equation}
  \mathbf{a}_{\text{aero}}\cdot\vhat = -\frac{D}{m}, \qquad
  \mathbf{a}_{\text{aero}}\cdot\ig   = \frac{L\cos\sigma}{m}, \qquad
  \mathbf{a}_{\text{aero}}\cdot\ip   = \frac{L\sin\sigma}{m} .
\end{equation}

\paragraph{Physical reading.}
$\sigma=0$ is wings-level: all the lift holds the vehicle up and none of it
turns. $\sigma=\pi/2$ is a knife-edge roll: all the lift turns and none of it
holds the vehicle up, so the vehicle falls as fast as a ballistic body while
turning as hard as it can. The descent phase of \code{run\_boost\_glide} uses
$\sigma=75^{\circ}$ for exactly this reason: $\cos 75^{\circ}=0.2588$, so it
keeps only a quarter of its lift supporting the vehicle and dives. The sign of
$\sigma$ chooses left or right; a sign error there mirrors the ground track
about the launch great circle and, on a zero-bank configuration, is completely
invisible.

The $\cos\gamma$ in the denominator of \eqref{eq:psigen2} means the same bank
angle turns the heading much faster in steep flight, and diverges at
$\gamma=\pm\pi/2$. This is again a coordinate singularity --- in vertical flight
``heading'' is undefined --- and it is guarded at $|\cos\gamma|<10^{-8}$,
raising \code{verticalFlight}.

\subsection{Thrust}
\label{sec:thrust}

Thrust acts at angle of attack $\alpha$ from the velocity vector, in the plane
set by the same bank angle $\sigma$ that orients the lift. That is a modelling
choice --- it says the vehicle's body axis, its lift vector and its thrust
vector all share one roll orientation --- and its payoff is that thrust and
lift add channel by channel with no extra bookkeeping:
\begin{equation}
  \mathbf{a}_{\text{thr}} \;=\; \frac{T}{m}\Bigl[
    \cos\alpha\,\vhat + \sin\alpha\bigl(\cos\sigma\,\ig
    + \sin\sigma\,\ip\bigr)\Bigr].
  \label{eq:athr}
\end{equation}
The three projections follow at once from orthonormality:
\begin{align}
  \mathbf{a}_{\text{thr}}\cdot\vhat &= \frac{T\cos\alpha}{m},
    \label{eq:tv}\\
  \mathbf{a}_{\text{thr}}\cdot\ig   &= \frac{T\sin\alpha\cos\sigma}{m},
    \label{eq:tg}\\
  \mathbf{a}_{\text{thr}}\cdot\ip   &= \frac{T\sin\alpha\sin\sigma}{m},
    \label{eq:tp}
\end{align}
and $\|\mathbf{a}_{\text{thr}}\| = T/m$ identically, since
$\cos^2\!\alpha + \sin^2\!\alpha(\cos^2\!\sigma+\sin^2\!\sigma) = 1$. All three
were verified symbolically, and numerically against the library at
$\alpha=0.13$, $\sigma=-0.62$ to a worst relative residual of
$1.6\times10^{-16}$; the norm identity holds to $1.2\times10^{-10}~\mathrm{N}$
on a $9.5\times10^{5}~\mathrm{N}$ motor, which is floating-point noise.

\paragraph{Physical reading.} At $\alpha=0$ the motor pushes purely along the
velocity: maximum speed gain, no turning. As $\alpha$ opens, $\cos\alpha$ costs
speed as $1-\alpha^2/2$ --- second order, hence cheap for small incidence ---
while $\sin\alpha$ buys turning as $\alpha$ --- first order, hence effective.
That asymmetry is the whole basis of gravity-turn ascent: a few degrees of
incidence steers the vehicle for almost no velocity penalty. The library's
measured steering loss on the shipped ballistic ascent,
$\int T(1-\cos\alpha)/m\,\dd t$, is $33.5~\mathrm{m/s}$ out of a
$6636~\mathrm{m/s}$ ideal --- half a percent.

Dropping $\cos\alpha$ from \eqref{eq:tv} is a small error that no reduction
test can see (Section~\ref{sec:phases}); swapping $\cos\sigma$ and $\sin\sigma$
between \eqref{eq:tg} and \eqref{eq:tp} rotates the whole thrust vector by
$90^{\circ}$ about the velocity and is a large one. The library's thrust
increment check exists to catch exactly these two, and when they are injected
as mutations it measures them at $9.8\times10^{-3}$ relative and
$30.0\,\%$ respectively.

\subsection{Mass}

Propellant leaves the vehicle at a positive rate $\dot m_p$, so
\begin{equation}
  \dot m \;=\; -\dot m_p .
  \label{eq:mdot}
\end{equation}
The sign convention is fixed at the interface: the propulsion model returns
$\dot m_p$ \emph{positive}, and the equations of motion apply the minus. Both
halves of that convention have to be written down, because each is defensible
alone and only the pair is consistent.

\paragraph{Why the mass in every acceleration denominator must be the state.}
Equations \eqref{eq:aaero} and \eqref{eq:athr} divide by $m$. In
\code{boost3DOF} that $m$ is $\mathbf{x}_7$ --- the integrated state --- and not
a field of the parameter struct. This is not a stylistic preference. A booster
burning $30{,}000~\mathrm{kg}$ of propellant from a $32{,}400~\mathrm{kg}$
liftoff mass ends at $2{,}400~\mathrm{kg}$: a factor of $13.5$. Freezing the
denominator at the liftoff value understates the end-of-burn acceleration by
that factor and costs $64\,\%$ of the delta-V. Freezing it at the burnout value
overstates the early acceleration by as much.

The trap is that the six \emph{unpowered} equations read the mass from
\code{veh.mass} instead, because they were written before a mass state existed.
A chain that carries $\mathbf{x}_7$ through an unpowered phase therefore holds
one physical quantity in two places, with nothing in the arithmetic forcing
them to agree --- a staging jump applied to the state alone leaves the flight
running silently at the old weight. Section~\ref{sec:phases} describes the
guard that closes this, and \code{docs/LESSONS\_LEARNED.md} records the general
form of the lesson.

\subsection{The assembled equations}
\label{sec:assembled}

Collecting \eqref{eq:rdot}, \eqref{eq:latlondot}, \eqref{eq:Vgen2}--%
\eqref{eq:psigen2}, \eqref{eq:gv}--\eqref{eq:gp}, \eqref{eq:aaero},
\eqref{eq:athr} and \eqref{eq:mdot}, and writing $a_L = L/m$, $a_D = D/m$,
$a_{T\parallel} = T\cos\alpha/m$, $a_{T\perp} = T\sin\alpha/m$, the
seven-state powered equations of motion over a non-rotating sphere are
\begin{align}
  \dot r        &= V\sin\gamma ,                                        \label{eq:f1}\\
  \dot \lambda  &= \frac{V\cos\gamma\sin\psi}{r\cos\phi} ,              \label{eq:f2}\\
  \dot \phi     &= \frac{V\cos\gamma\cos\psi}{r} ,                      \label{eq:f3}\\
  \dot V        &= a_{T\parallel} - a_D - g_r\sin\gamma
                   + g_\phi\cos\gamma\cos\psi ,                         \label{eq:f4}\\
  \dot \gamma   &= \frac{(a_{T\perp}+a_L)\cos\sigma}{V}
                   - \Bigl(\frac{g_r}{V} - \frac{V}{r}\Bigr)\cos\gamma
                   - \frac{g_\phi\sin\gamma\cos\psi}{V} ,               \label{eq:f5}\\
  \dot \psi     &= \frac{(a_{T\perp}+a_L)\sin\sigma}{V\cos\gamma}
                   + \frac{V\cos\gamma\sin\psi\tan\phi}{r}
                   - \frac{g_\phi\sin\psi}{V\cos\gamma} ,               \label{eq:f6}\\
  \dot m        &= -\dot m_p .                                          \label{eq:f7}
\end{align}
Setting $T=0$ (equivalently $a_{T\parallel}=a_{T\perp}=0$) and dropping
\eqref{eq:f7} gives \code{glide3DOF} exactly. These are
\cite{Vinh1980}~Eqs.~(2.28)--(2.33) in this library's notation.

Two of the groupings in \eqref{eq:f5} are worth naming, because they are the
places where the code's algebra departs cosmetically from the derivation
above. The derivation produced $\mathbf{a}\cdot\ig/V + (V/r)\cos\gamma$; the
code writes the gravity part and the transport part as one bracket,
$-(g_r/V - V/r)\cos\gamma$. They are the same expression. Likewise
\eqref{eq:f6}'s middle term is the transport term of \eqref{eq:transport}
divided by $\cos\gamma$, which cancels one power of the cosine.

\subsection{Centrifugal relief and the turning term}
\label{sec:relief}

\paragraph{$\bigl(g_r/V - V/r\bigr)\cos\gamma$: centrifugal relief.}
In level flight ($\gamma=0$, no lift) this term is the whole of
$\dot\gamma$. It vanishes when
\begin{equation}
  \frac{g_r}{V} = \frac{V}{r}
  \quad\Longleftrightarrow\quad
  V^2 = g_r r \;=\; \frac{\mu}{r}
  \quad\Longleftrightarrow\quad
  V = V_c \equiv \sqrt{\mu/r} ,
  \label{eq:circular}
\end{equation}
the local circular speed, $V_c = 7905.4~\mathrm{m/s}$ at the reference sphere.
The reading is exact and physical: at circular speed the required centripetal
acceleration $V^2/r$ is supplied entirely by gravity and the path curves at
precisely the rate needed to follow the sphere, so no lift is required. Below
$V_c$ gravity wins and the path bends down; above $V_c$ the vehicle is in
effect on a rising orbit and the path bends up even with no lift at all.

For a hypersonic glider the practical consequence is that only the
\emph{deficit} $1 - V^2/V_c^2$ has to be carried by lift. At $V=6000$~m/s that
deficit is $0.424$: the vehicle weighs, aerodynamically, less than half of what
it does at rest, which is why an $L/D$ of $2.5$ is enough to fly halfway around
the planet (Section~\ref{sec:saenger}). Get the sign of this term wrong and the
vehicle behaves as though gravity were \emph{stronger} at high speed --- the
glide collapses, and the equilibrium-glide check of
Section~\ref{sec:saenger} fails by tens of percent rather than the measured
$1.5\,\%$.

\paragraph{$V\cos\gamma\sin\psi\tan\phi/r$: the turning term.}
This term appears in $\dot\psi$ with no force behind it. It is pure geometry:
holding a great-circle course over a sphere requires the heading to change
continuously, because the meridians the heading is measured against are
themselves converging. Its structure says so. It is proportional to
$\sin\psi$, so it vanishes for a due-north or due-south leg --- a meridian
\emph{is} a great circle of constant bearing. It is proportional to $\tan\phi$,
so it vanishes on the equator --- also a great circle of constant bearing ---
and grows without bound at the pole, where all meridians meet. And it is
proportional to the horizontal ground speed $V\cos\gamma$ divided by $r$, which
is the angular rate at which the vehicle is crossing those meridians.

An eastbound flight in the northern hemisphere ($\sin\psi>0$, $\tan\phi>0$) has
$\dot\psi>0$: its bearing rotates clockwise, from east toward south. That is
the familiar great-circle behaviour --- the route from Los Angeles to New York
departs on a bearing of $65.87^{\circ}$ and arrives on one of $93.84^{\circ}$
($=273.84^{\circ}-180^{\circ}$), swinging nearly $28^{\circ}$ to the south
along the way. Delete the term and a zero-bank flight no longer follows a great
circle at all; the great-circle range checks of Section~\ref{sec:gc} are what
notice.

\subsection{Rotating Earth}
\label{sec:rotating}

Everything above is written in the rotating frame, where $V$ is the
planet-relative speed. Two apparent accelerations are then required, and they
are gated on \code{env.omegaE} so that setting it to zero recovers the
non-rotating case \emph{exactly} --- not approximately, since the whole block
is skipped:
\begin{equation}
  \mathbf{a}_{\text{rot}} \;=\;
  \underbrace{-2\,\bOmega\times\mathbf{V}}_{\text{Coriolis}}
  \;\underbrace{-\,\bOmega\times(\bOmega\times\mathbf{r})}_{\text{centrifugal}},
  \qquad \bOmega = \omega\,\zhat .
  \label{eq:arot}
\end{equation}

\paragraph{Centrifugal.}
$-\bOmega\times(\bOmega\times\mathbf{r})$ points directly away from the spin
axis with magnitude $\omega^2 r\cos\phi$. Resolving the outward-from-axis
direction on the horizon basis as $\cos\phi\,\uhat - \sin\phi\,\nhat$ and
projecting:
\begin{align}
  \mathbf{a}_{\text{cen}}\cdot\vhat
    &= \omega^2 r\cos\phi\,\bigl(\sin\gamma\cos\phi
       - \cos\gamma\sin\phi\cos\psi\bigr) ,                      \label{eq:cenv}\\
  \mathbf{a}_{\text{cen}}\cdot\ig
    &= \omega^2 r\cos\phi\,\bigl(\cos\gamma\cos\phi
       + \sin\gamma\cos\psi\sin\phi\bigr) ,                      \label{eq:ceng}\\
  \mathbf{a}_{\text{cen}}\cdot\ip
    &= \omega^2 r\sin\phi\cos\phi\sin\psi ,                      \label{eq:cenp}
\end{align}
which enter $\dot V$, $\dot\gamma$ and $\dot\psi$ divided by $1$, $V$ and
$V\cos\gamma$ respectively, exactly as \eqref{eq:Vgen2}--\eqref{eq:psigen2}
require. All three match the source.

\paragraph{Coriolis.}
$-2\bOmega\times\mathbf{V}$ is perpendicular to $\mathbf{V}$ by construction,
so it appears in \emph{neither} $\dot V$ nor the kinematics. Its two surviving
projections reduce to
\begin{align}
  \frac{\mathbf{a}_{\text{cor}}\cdot\ig}{V}
    &= 2\omega\cos\phi\sin\psi ,                                  \label{eq:corg}\\
  \frac{\mathbf{a}_{\text{cor}}\cdot\ip}{V\cos\gamma}
    &= 2\omega\bigl(\sin\phi - \cos\phi\cos\psi\tan\gamma\bigr) .  \label{eq:corp}
\end{align}
Both are independent of $V$ --- the speed cancels between the force and the
$1/V$ of \eqref{eq:gamgen2} --- so Coriolis is a pure \emph{rate} of turn,
$O(2\omega) = 1.5\times10^{-4}~\mathrm{rad/s}$, or about half a degree per
minute. Over a 27-minute intercontinental flight that is up to $14^{\circ}$ of
accumulated heading, reduced by the geometric factors in
\eqref{eq:corg}--\eqref{eq:corp}. It is why an azimuth solved on a
non-rotating Earth misses (Section~\ref{sec:targeting}).

Reading \eqref{eq:corg}: an eastbound vehicle ($\sin\psi>0$) is pushed
\emph{up}, a westbound one down --- the E\"otv\"os effect, and it is largest at
the equator. Reading the first part of \eqref{eq:corp}: $2\omega\sin\phi$ is
the classical Coriolis parameter, deflecting to the right in the northern
hemisphere and to the left in the southern, vanishing on the equator. The
second part carries a $\tan\gamma$ and so blows up in vertical flight; it is
covered by the same $|\cos\gamma|$ guard.

Note that \eqref{eq:cenv}--\eqref{eq:corp} are all multiplied out and added to
the non-rotating right-hand side inside an \code{if om \textasciitilde= 0}
block. With the shipped \code{env.omegaE = 0} they are \textbf{zero by
default}: every trajectory in the library's documentation, and every entry
script except by explicit override, is flown on a non-rotating sphere.

\subsection{Singularities and guards}

Three states make the equations singular and all three are guarded explicitly,
raising a named error rather than returning \code{NaN}:

\begin{center}
\begin{tabular}{lll}
\toprule
Condition & Singular quantity & Identifier suffix \\
\midrule
$|\cos\phi| < 10^{-8}$   & $\dot\lambda$, Eq.~\eqref{eq:f2}            & \code{polarSingularity} \\
$|\cos\gamma| < 10^{-8}$ & $\dot\psi$, Eq.~\eqref{eq:f6}               & \code{verticalFlight} \\
$V < 1~\mathrm{m/s}$     & $\dot\gamma$, $\dot\psi$, Eqs.~\eqref{eq:f5}--\eqref{eq:f6} & \code{zeroSpeed} \\
$m \le 0$ (powered only) & every $1/m$, Eqs.~\eqref{eq:aaero}, \eqref{eq:athr} & \code{massDepleted} \\
\bottomrule
\end{tabular}
\end{center}

The last is not a formality. A negative mass state does not merely scale the
accelerations, it \emph{inverts} them --- thrust becomes a retarding force and
drag becomes a propulsive one --- so the guard is written $m\le 0$ and catches
the boundary case too. The $V<1$ guard is why the entry scripts start their
ascent at a nonzero speed (10~m/s in \code{run\_ballistic}); a launch from
literal rest is outside the model, not merely awkward for it.

%=====================================================================
\section{The injected models}
\label{sec:models}

The equations of motion never name a model. They call four handles carried on
an environment struct, each with a fixed signature so that a higher-fidelity
replacement is a one-line change in an entry script:

\begin{center}
\begin{tabular}{lll}
\toprule
Family & Signature & Shipped model \\
\midrule
atmosphere   & \code{[rho,P,T,a] = f(h)}          & \code{expAtmos} \\
gravity      & \code{[gr,gLat]   = f(r,lat)}      & \code{sphereGrav} \\
aerodynamics & \code{[CL,CD]     = f(alpha,mach,veh)} & \code{constLD} \\
propulsion   & \code{[T,mdot]    = f(t,P,veh)}    & \code{constThrust} \\
\bottomrule
\end{tabular}
\end{center}

\subsection{Exponential isothermal atmosphere}

\begin{equation}
  \rho(h) = \rho_0\,e^{-h/H}, \qquad
  T(h) = T_0, \qquad
  P = \rho R T_0, \qquad
  a_s = \sqrt{\kappa R T_0},
  \label{eq:atmos}
\end{equation}
with $\rho_0 = 1.225~\mathrm{kg/m^3}$, $H = 7200~\mathrm{m}$,
$T_0 = 250~\mathrm{K}$, $R = 287.0528~\mathrm{J/kg/K}$, $\kappa = 1.4$. Those
give $P(0) = 87{,}909.92~\mathrm{Pa}$ and $a_s = 316.97~\mathrm{m/s}$
everywhere --- the sound speed is altitude-independent because the temperature
is.

\paragraph{This is not exactly hydrostatic, and that is an accepted Phase~1
choice.} Hydrostatic equilibrium in an isothermal gas requires
$\dd P/\dd r = -\rho\,g(r)$, i.e.\ $\dd\ln\rho/\dd r = -g(r)/(RT_0)$. Two
separate approximations stand between that and \eqref{eq:atmos}:

\begin{enumerate}
\item \emph{The scale height is not the hydrostatic one.} With constant gravity
the solution of the hydrostatic equation is exponential with
$H_{\text{hyd}} = RT_0/g_0 = 7317.81~\mathrm{m}$. The library uses
$7200~\mathrm{m}$, which is $117.81~\mathrm{m}$ ($1.61\,\%$) smaller. The
$7200~\mathrm{m}$ figure is an empirical fit to the real low-altitude density
profile, not a derived quantity, and it is the density profile --- not the
pressure --- that the trajectory actually feels.
\item \emph{Gravity is not constant.} Under $g(r)=\mu/r^2$ the exact isothermal
hydrostatic profile is not an exponential at all but
$\rho/\rho_0 = \exp\bigl[(\mu/RT_0)(1/r - 1/r_E)\bigr]$, whose local scale
height grows with altitude ($7555~\mathrm{m}$ at 100~km against
$7318~\mathrm{m}$ at the surface). Against that exact profile the shipped
exponential is low by $8.6\,\%$ at 30~km, $20\,\%$ at 60~km and $36\,\%$ at
100~km.
\end{enumerate}

Neither discrepancy is hidden and neither is a bug: the model's stated purpose
is to be the atmosphere for which the entry problem has closed-form solutions
(Sections~\ref{sec:saenger} and \ref{sec:allen}), and Allen--Eggers needs an
exponential \emph{density} profile and nothing else --- temperature and
pressure never enter it. What the mismatch does mean is that $P$ from
\eqref{eq:atmos} is a weaker quantity than $\rho$, and the one place $P$ is
consumed is the thrust back-pressure debit below, where an error of a few
percent in $P$ moves the delivered thrust by a few tenths of a percent.

\paragraph{The trap.} $H$ appears both in the model and in every closed-form
entry solution derived from the model, so it \emph{cancels} out of those
comparisons: inflating $H$ by $15\,\%$ moves the Allen--Eggers agreement from
$+10.50\,\%$ to $+10.44\,\%$. No tolerance recovers a sensitivity that is
identically zero. $H$ is therefore pinned by exact equality at its source, in
\code{test\_missileConst}, and that pin is the only thing protecting it.

\subsection{Point-mass gravity}

\begin{equation}
  g_r = \frac{\mu}{r^2} \;>\; 0 , \qquad g_\phi \equiv 0 ,
  \label{eq:grav}
\end{equation}
with $\mu = 3.986004418\times10^{14}~\mathrm{m^3/s^2}$ (WGS-84), giving
$9.798285~\mathrm{m/s^2}$ at the reference sphere --- note this is
\emph{not} $g_0 = 9.80665$, which is a defined standard used only for
converting $I_{sp}$ and for reporting load factors in $g$.

The interface carries $g_\phi$ even though it is identically zero because a
scalar cannot represent $J_2$. Committing to a scalar now would force an edit
to \eqref{eq:f4}--\eqref{eq:f6} the day an oblate model arrives; carrying the
zero costs three multiplications by zero per derivative call and buys a drop-in
replacement. The cost, stated plainly, is that the $g_\phi$ terms are
\emph{dead code under the shipped configuration} and are unexercised by the
test suite --- see the caveat in Section~\ref{sec:gravity}.

\subsection{Constant lift coefficient, constant lift-to-drag ratio}

\begin{equation}
  C_L = \text{const}, \qquad C_D = \frac{C_L}{(L/D)} .
  \label{eq:aero}
\end{equation}
$\alpha$ and Mach are accepted and ignored. Two consequences follow and both
matter downstream.

First, there is deliberately \emph{no} $C_D$ field on any vehicle struct: drag
is derived from the pair, so it cannot fall out of sync with lift. Expressing
``almost no lift, plenty of drag'' for a re-entry body therefore requires a
small $C_L$ over a small $L/D$ --- \code{vehicle\_bm} uses
$C_L=0.005$, $L/D=0.02$, giving $C_D = 0.25$ --- and $C_L$ cannot be set to
zero without taking $C_D$ to zero with it.

Second, because $\alpha$ does not reach the aerodynamics, the \emph{only} live
aerodynamic control in the library is the bank angle $\sigma$. The angle of
attack still reaches the \emph{thrust} projections \eqref{eq:tv}--\eqref{eq:tp},
so it is a genuine control during boost and an inert one during glide. Anyone
reading a glide phase's $\alpha$ history and expecting it to matter will be
disappointed; that changes the day a table-driven $C_L(\alpha,M)$ is installed,
at which point nothing else has to change.

\subsection{Choked-nozzle thrust with back-pressure correction}

\begin{align}
  T &= \max\bigl(T_{\text{vac}} - A_e P,\; 0\bigr) ,          \label{eq:thrust}\\
  \dot m_p &= \frac{T_{\text{vac}}}{I_{sp}\,g_0}\;\mathbf{1}\{T>0\} .
              \label{eq:mflow}
\end{align}

\paragraph{Why $\dot m_p$ is held fixed while $T$ varies.}
This looks inconsistent on first reading and is in fact the physically correct
pairing. The nozzle throat is choked and the chamber pressure is fixed, so the
mass flow through the throat is set entirely by conditions \emph{upstream} of
the throat. Information cannot travel upstream through a sonic throat, so
ambient back pressure --- which acts downstream, on the exit plane --- cannot
change the flow. What it changes is the delivered thrust, through the pressure
term $A_e(P_e - P)$ acting on the exit area. Hence: flow constant, thrust
altitude-dependent.

The same argument fixes which $I_{sp}$ may be used. The identity
$\dot m_p = T/(I_{sp}g_0)$ holds only when thrust and specific impulse are
quoted at the \emph{same} ambient condition; \eqref{eq:mflow} uses the vacuum
pair, so it is consistent, and the flow it computes is the flow at every
altitude. Substituting a sea-level $I_{sp}$ against a vacuum thrust would
overstate the flow by the sea-level debit and burn the motor out early --- a
units error that presents as a trajectory error.

\paragraph{The clamp zeroes the flow too.}
The linear debit $A_e P$ is only valid while it is small. Driven far enough it
would return a negative thrust, which the $\max$ clamps to zero; at that point
the model has left its domain, so the motor is reported as \emph{not running},
and $\dot m_p$ goes to zero with it. A clamped engine that kept consuming
propellant would corrupt the mass history silently and fire the burnout event
early. The two outputs are always zero together.

\paragraph{Numbers for the shipped placeholder booster.}
$T_{\text{vac}}=950{,}000~\mathrm{N}$, $I_{sp}=260~\mathrm{s}$,
$A_e=1.25~\mathrm{m^2}$, $m_{\text{prop}}=30{,}000~\mathrm{kg}$,
$m_{\text{dry}}=1{,}500~\mathrm{kg}$ over a $900~\mathrm{kg}$ payload. Then
\begin{align*}
  \dot m_p &= \frac{950000}{260\times 9.80665} = 372.5886162804~\mathrm{kg/s},
  & t_{\text{burn}} &= \frac{30000}{372.5886} = 80.517758~\mathrm{s}, \\
  T(\text{sea level}) &= 840{,}112.6~\mathrm{N},
  & \text{debit} &= 109{,}887~\mathrm{N} = 11.57\,\% ,\\
  \text{vacuum } T/W &= \frac{950{,}000}{32400\,g_0} = 2.9899 ,
  & \text{liftoff } T/W &= \frac{840{,}112.6}{32400\,g_0} = 2.6441 .
\end{align*}
All were re-computed for this note and agree with the values pinned in the
test suite.

\paragraph{Both thrust-to-weight figures, and which one the vehicle feels.}
Quoting one without the other is the standard way to be off by a tenth at the
pad. $T_{\text{vac}}$ is never delivered at sea level --- the same
$A_eP$ debit computed two lines above takes it to $840{,}112.6~\mathrm{N}$ ---
so $2.9899$ is the data-sheet number and $2.6441$ is what actually lifts the
stack. The difference is not cosmetic: net of weight, the initial vertical
acceleration is $1.6441\,g$ rather than $1.9899\,g$, a $17\,\%$ shortfall on
the quantity that has to clear the tower. The vacuum figure is the right one
to compare against a motor specification; the delivered figure is the right one
to put in a trajectory.

%=====================================================================
\section{Phase machinery}
\label{sec:phases}

\subsection{The seven-state convention}

A trajectory is a sequence of phases, each an ODE integration with its own
right-hand side, its own guidance law, its own terminal event and its own time
span. \code{coorbital.prop.phaseRun} runs them in order, handing the terminal
state of one to the next and concatenating the results into a single
\code{traj.x}. Because that concatenation produces one rectangular matrix,
\emph{every phase in one call must share one state dimension}.

The consequence is that a chain containing a powered phase runs seven-state
throughout, and the six-state unpowered equations are lifted to match by
\code{coorbital.eom.massConstant}, which appends $\dot m = 0$. Unpowered phases
therefore carry an inert mass. This was chosen over the alternative --- a
mapping layer translating between a six-state and a seven-state vector at each
junction --- for a specific reason: a mapping at a seam is a place for state to
be silently lost or reordered, and a uniform vector has no such place. There is
nothing to get wrong at a seam because nothing happens at a seam.

\paragraph{The cost, which is real.} The lifted equations divide by
\code{veh.mass} and never read $\mathbf{x}_7$. One physical quantity is now
held in two places and nothing in the arithmetic makes them agree. A chain that
jettisons a booster by dropping $\mathbf{x}_7$, while continuing to hand the
same vehicle struct to the next phase, flies the rest of the flight at the old
weight --- with a mass history that says otherwise, no divergence, no warning,
and an entirely plausible trajectory. Measured on this library: deleting the
separation link left the flown trajectory \emph{bit-identical} and changed only
the recorded mass.

\code{massConstant} therefore refuses to run when the carried mass and the
flown mass disagree by more than $10^{-9}$ relative, floored at $10^{-9}$
absolute. Two design rules are attached to that guard and both were learned by
getting them wrong first:
\begin{itemize}
\item \textbf{Raise; do not reconcile.} Injecting $\mathbf{x}_7$ into a copy of
the vehicle struct is the obvious repair and is worse than the bug, because it
\emph{looks} repaired: separation changes $S$, $C_L$ and $L/D$ as well as the
mass, and silencing the one symptom that is visible converts a detectable fault
into an undetectable one. The structural answer is a per-phase vehicle bound
inside each phase's EOM closure, which is what all four entry scripts do.
\item \textbf{Floor the tolerance.} A purely relative budget collapses as the
mass approaches zero. At $0.25~\mathrm{kg}$ it would be
$2.5\times10^{-10}~\mathrm{kg}$, tighter than the
$3.4\times10^{-10}~\mathrm{kg}$ residual an ODE burnout-event solve already
leaves on a 900~kg stack --- the guard would then reject ordinary roundoff.
\end{itemize}

\subsection{Stage separation is a genuine state discontinuity}

A phase may carry an optional link, $\mathbf{x}^+ = \ell(\mathbf{x}^-)$, applied
to its terminal state to form the next phase's initial state. Stage separation
is one line:
\begin{equation}
  \ell(\mathbf{x}) = \begin{bmatrix} \mathbf{x}_{1:6} \\
                     \mathbf{x}_7 - m_{\text{dry}} \end{bmatrix} .
  \label{eq:link}
\end{equation}
This is not a numerical convenience standing in for a fast transient. The
booster physically leaves, and the state is genuinely discontinuous in
$\mathbf{x}_7$ at that instant, while being continuous in the six flight
states. A test asserting continuity in all seven components would be asserting
a falsehood; the library asserts a jump of exactly $m_{\text{dry}}$ in
component~7 and continuity in components 1--6 against an independent
higher-order re-integration across the seam.

The link is applied \emph{before} the junction is recorded, so
\code{traj.junction(k).x} is the state as handed to the next phase --- the far
side of the jump. The near side is the last trajectory row labelled phase $k$,
at the same instant on the cumulative clock.

\subsection{ODE events and the one-sided \code{direction}}

Each phase terminates on a MATLAB ODE event function returning
$(\text{value},\text{isterminal},\text{direction})$. The three shipped events
are
\begin{center}
\begin{tabular}{lll}
\toprule
Event & Value & \code{direction} \\
\midrule
\code{eventAltitude(t,x,hStop)}   & $(r - r_E) - h_{\text{stop}}$ & $-1$ \\
\code{eventBurnout(t,x,mBurnout)} & $\mathbf{x}_7 - m_{\text{burnout}}$ & $-1$ \\
\code{eventApogee(t,x)}           & $\gamma$ & $-1$ \\
\bottomrule
\end{tabular}
\end{center}
All are terminal and all are \emph{one-sided}. The \code{direction} field is
the load-bearing part and \code{-1} means ``locate only crossings at which the
value is decreasing''.

Consider what happens without it. A lofted trajectory launched below a trigger
altitude climbs \emph{through} that altitude on the way up, coasts to apogee
far above it, and descends back through it. A two-sided event stops the run on
the ascending crossing: the flight ends near the launch point with a peak
altitude barely above the trigger and a \emph{positive} terminal flight-path
angle. Nothing about the result is obviously wrong --- an altitude event fired
at the requested altitude --- and the whole coast and descent are simply
missing. The one-sided form makes that impossible. The library tests it
behaviourally, by flying an arc that must cross the trigger twice and asserting
that the peak altitude exceeds it by tens of kilometres.

Apogee deserves a note of its own. Terminating on $\gamma$ crossing zero
downward makes apogee an \emph{exactly solved} event rather than the largest
sample on an adaptive grid, so apogee time and altitude are good to the event
solver's tolerance rather than to a step size. That is why
\code{run\_ballistic} splits its coast in two at apogee even though the physics
does not require a split; it also leaves the final phase monotonically
descending, so the one-sided impact event cannot be reached from the wrong
side.

\subsection{Guidance}

Two laws are shipped. \code{prescribed} interpolates $(\alpha,\sigma)$ linearly
on a time grid and clamps outside it; it ignores the state.
\code{pitchProgram} is the first law that reads the state: it interpolates a
commanded pitch \emph{attitude} $\theta(t)$ and returns
\begin{equation}
  \alpha = \theta(t) - \gamma, \qquad \gamma = \mathbf{x}_5 ,
  \label{eq:pitchprog}
\end{equation}
optionally clamped to $|\alpha|\le\alpha_{\max}$. This is the classical
prescribed-pitch ascent law: the vehicle is flown to an attitude schedule and
the angle of attack is whatever gap opens between that attitude and the
velocity vector. The clamp is not optional in practice --- a $90^{\circ}$
commanded attitude at liftoff, when $\gamma$ is still near zero, would hand the
equations of motion a physically absurd incidence.

One convention must be internalised: each phase's guidance is evaluated at that
phase's \emph{own} tspan-relative time, not at cumulative flight time. A glide
schedule written against $[0,t_{\max}]$ therefore starts at the handoff from
boost, not at liftoff. With constant schedules the distinction is invisible;
with a time-varying bank it moves the range by 30\,\%.

%=====================================================================
\section{Analytic references}
\label{sec:analytic}

Seven closed forms stand behind the library. Each is derived here and each is
reported with the agreement measured on the committed code. Two properties of
this list are worth stating before the list itself. First, several of these
references are computed from the same constants as the model, so they are
structurally blind to those constants: the constant appears on both sides and
cancels, and no tolerance recovers the sensitivity. Where that happens it is
said so explicitly. Second, no single check covers the model --- the
equilibrium-glide reference contains no drag term and is blind to $C_D$; the
Allen--Eggers reference is zero-lift and is blind to $C_L$. They are
complementary by construction.

\subsection{Vacuum specific energy}
\label{sec:energy}

With no atmosphere, no lift and no rotation, the only force is point-mass
gravity, which is conservative. The specific orbital energy
\begin{equation}
  \varepsilon \;=\; \frac{V^2}{2} - \frac{\mu}{r}
  \label{eq:energy}
\end{equation}
must therefore be constant. Directly, using \eqref{eq:f1} and \eqref{eq:f4}
with $g_r=\mu/r^2$, $g_\phi=0$, $a_D = a_T = 0$:
\begin{equation}
  \dot\varepsilon = V\dot V + \frac{\mu}{r^2}\dot r
  = V\Bigl(-\frac{\mu}{r^2}\sin\gamma\Bigr) + \frac{\mu}{r^2}\,V\sin\gamma
  = 0 .
\end{equation}
The cancellation is exact and involves only two terms, which is precisely what
makes this the sharpest single instrument in the suite: \emph{any} sign error
in the radial gravity term, or in $\dot r$, or in the relative sign of the two,
destroys it immediately. It is also cheap --- no second propagator is needed to
produce the reference.

\textbf{Measured:} $1.281\times10^{-13}$ relative drift over a 300~s lofted
arc, against a $10^{-8}$ budget.

\subsection{The rotating-frame Jacobi integral, and what it cannot see}
\label{sec:jacobi}

With the Earth rate on, energy is no longer conserved --- the frame is not
inertial --- but the Jacobi integral is:
\begin{equation}
  J \;=\; \frac{V^2}{2} - \frac{\mu}{r}
        - \frac{1}{2}\,\omega^2 r^2\cos^2\!\phi .
  \label{eq:jacobi}
\end{equation}
The third term is the centrifugal potential, $-\tfrac12\omega^2 d^2$ with
$d = r\cos\phi$ the distance from the spin axis. Differentiating,
\begin{equation}
  \frac{\dd J}{\dd t} \;=\; V\dot V + \frac{\mu}{r^2}\dot r
   - \omega^2 r\cos^2\!\phi\;\dot r
   + \omega^2 r^2\cos\phi\sin\phi\;\dot\phi .
  \label{eq:djdt}
\end{equation}
Substituting $\dot V$ from \eqref{eq:f4} with the centrifugal term
\eqref{eq:cenv}, $\dot r$ from \eqref{eq:f1} and $\dot\phi$ from
\eqref{eq:f3}, the four terms cancel in pairs:
\begin{align}
  V\dot V &= -\frac{\mu V\sin\gamma}{r^2}
             + \omega^2 r V\cos\phi\bigl(\sin\gamma\cos\phi
             - \cos\gamma\sin\phi\cos\psi\bigr), \nonumber\\
  \frac{\mu}{r^2}\dot r &= +\frac{\mu V\sin\gamma}{r^2}, \nonumber\\
  -\omega^2 r\cos^2\!\phi\,\dot r &= -\omega^2 r V\cos^2\!\phi\sin\gamma, \nonumber\\
  \omega^2 r^2\cos\phi\sin\phi\,\dot\phi
    &= +\omega^2 r V\cos\phi\sin\phi\cos\gamma\cos\psi ,
\end{align}
and $\dd J/\dd t = 0$ identically. This was confirmed symbolically.

\paragraph{The integral is structurally blind to a Coriolis sign error.}
Look at \eqref{eq:djdt} again. It contains $\dot V$, $\dot r$ and $\dot\phi$
and \emph{nothing else}: there is no $\dot\gamma$ and no $\dot\psi$ anywhere in
it, because $J$ does not depend on $\gamma$ or $\psi$. But the Coriolis
acceleration is workless --- $\mathbf{a}_{\text{cor}}\cdot\vhat = 0$ by
construction, verified numerically at $-1.4\times10^{-17}~\mathrm{m/s^2}$ ---
so it appears in \emph{neither} $\dot V$ nor the kinematic equations. It lives
only in $\dot\gamma$ and $\dot\psi$. The two facts compose: the Jacobi integral
cannot see the Coriolis terms at all.

This is not a hypothetical. Flipping the sign of both $2\omega$ terms by hand
and re-propagating the same arc:

\begin{center}
\begin{tabular}{lrr}
\toprule
& Correct Coriolis & Sign flipped \\
\midrule
Jacobi drift over 300~s   & $1.107\times10^{-13}$ & $1.468\times10^{-13}$ \\
Terminal radius change    & --- & $47{,}273~\mathrm{m}$ \\
Terminal latitude change  & --- & $2.940\times10^{-3}~\mathrm{rad}$ \\
\bottomrule
\end{tabular}
\end{center}

The conservation law is satisfied to the same $10^{-13}$ in both cases while
the trajectory has moved 47~km. A test built only on the Jacobi integral would
report green on a model with the Coriolis force pointing the wrong way. The
library closes the gap with an instantaneous spot check that recomputes
\eqref{eq:corg}--\eqref{eq:corp} directly at an off-equator, non-cardinal
state; for this note the same terms were rebuilt from explicit vector cross
products and agree to $1.1\times10^{-13}$ or better.

The general form of the lesson --- ask of any clean-looking test what would
have to be wrong for the number to move --- is recorded in
\code{docs/LESSONS\_LEARNED.md}.

\subsection{S\"anger equilibrium glide}
\label{sec:saenger}

\paragraph{The equilibrium condition.}
Set $\dot\gamma = 0$ in \eqref{eq:f5} with $\sigma=0$, no thrust and
$g_\phi=0$:
\begin{equation}
  \frac{L}{mV} = \Bigl(\frac{g_r}{V} - \frac{V}{r}\Bigr)\cos\gamma
  \quad\Longleftrightarrow\quad
  \tfrac12\rho V^2 S C_L = m\Bigl(g_r - \frac{V^2}{r}\Bigr)\cos\gamma .
  \label{eq:eqglide}
\end{equation}
Lift balances weight \emph{less centrifugal relief}, both projected by
$\cos\gamma$. Solving for the speed gives the closed form the library tests
against pointwise,
\begin{equation}
  V_{\text{eq}}^2(r) \;=\;
  \frac{m\,g_r(r)\cos\gamma}
       {\tfrac12\rho(r) S C_L + m\cos\gamma/r} .
  \label{eq:veq}
\end{equation}

\paragraph{The range: the approximations, stated before the integral.}
Equation~\eqref{eq:veq} is a pointwise condition and holds at whatever radius
it is evaluated at. The closed-form range that follows from it does not: three
approximations stand between the two, and they are written down here rather
than introduced silently as the integration proceeds.

\begin{enumerate}
\item \emph{Shallow flight.} $\gamma$ is small, so $\cos\gamma\to 1$ and the
component of gravity along the track is negligible beside drag.
\item \emph{Near-constant altitude.} The glide descends through a small
fraction of a planetary radius, so $r\simeq r_E$ over the whole arc. This is
the assumption that does the work, because it converts three varying
quantities into constants:
\begin{equation}
  r \simeq r_E , \qquad
  g \;\equiv\; g_r(r_E) = \frac{\mu}{r_E^2} = 9.798285~\mathrm{m/s^2},
  \qquad
  V_c^2 \;\equiv\; g\,r_E = \frac{\mu}{r_E},
  \quad V_c = 7905.4~\mathrm{m/s} .
  \label{eq:sangerconst}
\end{equation}
Note that $g$ is defined here, and defined \emph{as a constant}: it is the
inward attraction evaluated once at the reference sphere, not the $g_r(r)$ of
\eqref{eq:eqglide}. Without \eqref{eq:sangerconst} the ratio $V_c^2/g = r$
varies along the path and the substitution three lines below is not a
substitution at all. The approximation is usable because the variation is
small: over the library's 45~km to 20~km test glide, $r$ changes by
$0.39\,\%$ and $g$ by $0.78\,\%$.
\item \emph{Fixed $L/D$}, which the shipped aerodynamic model \eqref{eq:aero}
supplies exactly, so this one costs nothing here.
\end{enumerate}

With \eqref{eq:sangerconst} in force, write \eqref{eq:eqglide} at
$\cos\gamma = 1$ as
\begin{equation}
  L = mg\Bigl(1 - \frac{V^2}{V_c^2}\Bigr), \qquad V_c^2 = g\,r_E ,
  \label{eq:liftdeficit}
\end{equation}
which is the centrifugal-relief statement of Section~\ref{sec:relief}: lift
carries only the deficit. Energy along the ground track $s$, with the gravity
component negligible for shallow $\gamma$, is
\begin{equation}
  m V \frac{\dd V}{\dd s} = -D = -\frac{L}{(L/D)}
  = -\frac{mg}{(L/D)}\Bigl(1 - \frac{V^2}{V_c^2}\Bigr) .
\end{equation}
Separating,
\begin{equation}
  \dd s = -\frac{(L/D)}{g}\,\frac{V\,\dd V}{1 - V^2/V_c^2} .
\end{equation}
Substitute $u = V^2/V_c^2$, so $V\,\dd V = \tfrac12 V_c^2\,\dd u$ and,
crucially, $V_c^2/g = r_E$ --- a constant, by \eqref{eq:sangerconst}, which is
the only reason it can come outside the integral:
\begin{equation}
  \dd s = -\frac{(L/D)}{g}\cdot\frac{V_c^2}{2}\cdot\frac{\dd u}{1-u}
        = -\frac{(L/D)\,r_E}{2}\cdot\frac{\dd u}{1-u} .
  \label{eq:dsdu}
\end{equation}
Integrating from $u_0 = (V/V_c)^2$ down to $u=0$:
\begin{equation}
  R \;=\; \frac{L}{D}\,\frac{r_E}{2}\,
          \ln\!\left[\frac{1}{1 - (V/V_c)^2}\right].
  \label{eq:sanger}
\end{equation}

\paragraph{The factor of one half is essential.} It comes from
$V\,\dd V = \tfrac12 V_c^2\,\dd u$ and from nowhere else. Omitting it doubles
the predicted range exactly --- an error that is easy to make and impossible to
detect by inspecting the formula's limits, since both versions diverge at
$V\to V_c$ and both vanish at $V\to 0$. It was checked two ways for this note:
symbolically in \code{sympy}, and by direct numerical quadrature of
\eqref{eq:dsdu} at $(V/V_c)^2 = 0.25$, $0.5$ and $0.9$, agreeing with
\eqref{eq:sanger} to $5.3\times10^{-14}$, $1.4\times10^{-14}$ and
$1.9\times10^{-13}$ relative respectively.

\paragraph{Physical reading, and what the divergence is telling you.}
The logarithm is the interesting part. Range is not proportional to kinetic
energy; it diverges as $V\to V_c$, because at circular speed the vehicle needs
no lift, generates no drag through the $L/D$ relation, and does not come down.
The \emph{marginal} return runs the same way, not the opposite way:
\begin{equation}
  \frac{\dd R}{\dd V} \;=\; \frac{L}{D}\,\frac{r_E\,V}{V_c^2 - V^2}
  \;\longrightarrow\; \infty \qquad\text{as}\quad V\to V_c^{-} ,
  \label{eq:dRdV}
\end{equation}
so every increment of speed taken nearer circular speed buys \emph{more} range
than the one before it, not less. At $L/D=2.5$ the same $100~\mathrm{m/s}$ is
worth $27~\mathrm{km}$ of range at $1000~\mathrm{m/s}$, $373~\mathrm{km}$ at
$6000~\mathrm{m/s}$, and $23{,}687~\mathrm{km}$ at $7800~\mathrm{m/s}$ --- more
than half the planet's circumference for the last of those, which is the tell.

That last figure is not a performance claim; it is the model announcing that it
has left its domain. Every assumption behind \eqref{eq:sanger} fails together
as $V\to V_c$. The required lift $mg(1-V^2/V_c^2)$ goes to zero, so the
equilibrium altitude \eqref{eq:veq} climbs out of the sensible atmosphere and
the near-constant-altitude assumption \eqref{eq:sangerconst} goes with it. With
no lift there is no $L/D$ drag either, so the fixed-$L/D$ atmospheric-glide
picture describes a vehicle that is no longer flying aerodynamically. And at
$V = V_c$ exactly the vehicle is in a circular orbit, which is not a glide at
all and has no finite range to predict. The formula is quantitative while lift
genuinely carries the weight deficit --- roughly $V/V_c \lesssim 0.9$ --- and
above that it is a statement about its own boundary.

With $L/D = 2.5$ and $V = 6000~\mathrm{m/s}$
($V/V_c = 0.759$), \eqref{eq:sanger} predicts $6842~\mathrm{km}$;
the shipped boost-glide chain flies $7663~\mathrm{km}$ from a
$5841~\mathrm{m/s}$ burnout, the excess coming from the initial altitude, the
skip phugoid and the boost's own downrange leg, none of which
\eqref{eq:sanger} contains.

\textbf{Measured:} the library does not test \eqref{eq:sanger} directly ---
it tests the pointwise condition \eqref{eq:veq}, which is the statement
\eqref{eq:sanger} is integrated from. A glide launched \emph{on} the
equilibrium manifold (entry speed $V_{\text{eq}}$, entry flight-path angle from
the tangency condition
$\sin\gamma_0 = -2(D/m)/(2g + \dd V_{\text{eq}}^2/\dd r)$, so no phugoid is
excited) tracks it from 45~km to 20~km with a worst-case relative speed
departure of $\mathbf{1.51\,\%}$ over the \emph{whole} arc, against a $3\,\%$
budget, while the density spans a factor of $32.2$ and the speed a factor of
$4.80$. The departure is largest at the low, slow end where the quasi-steady
assumption gives way. For calibration: a $20\,\%$ error in $C_L$ drives the
same measurement to $11.52\,\%$.

\subsection{Allen--Eggers ballistic entry}
\label{sec:allen}

\paragraph{Setup.} A zero-lift body enters steeply. Two approximations: the
flight-path angle is frozen at its entry value $\gamma_e$, and gravity is
dropped from the speed equation, leaving drag alone. Then from \eqref{eq:f4}
and \eqref{eq:f1},
\begin{equation}
  \frac{\dd V}{\dd t} = -\frac{\rho V^2}{2\beta}, \qquad
  \frac{\dd h}{\dd t} = V\sin\gamma_e, \qquad
  \beta \equiv \frac{m}{S C_D},
\end{equation}
so, dividing,
\begin{equation}
  \frac{\dd V}{V} = -\frac{\rho(h)\,\dd h}{2\beta\sin\gamma_e} .
\end{equation}
With $\rho = \rho_0 e^{-h/H}$ and integrating from $h=\infty$ (where
$V = V_e$) down to $h$:
\begin{equation}
  \ln\frac{V}{V_e}
  = -\frac{1}{2\beta\sin\gamma_e}\int_{\infty}^{h}\rho_0 e^{-h'/H}\dd h'
  = \frac{H\rho(h)}{2\beta\sin\gamma_e} ,
\end{equation}
and since $\sin\gamma_e<0$ on entry this is negative, as it must be:
\begin{equation}
  V(\rho) = V_e\,\exp\!\left[-\frac{\rho H}{2\beta|\sin\gamma_e|}\right].
  \label{eq:aeV}
\end{equation}

\paragraph{Peak deceleration.} The aerodynamic deceleration is
$a = \rho V^2/(2\beta)$. Substituting \eqref{eq:aeV} and writing
$x \equiv \rho H/(\beta|\sin\gamma_e|)$,
\begin{equation}
  a(x) \;=\; \frac{V_e^2|\sin\gamma_e|}{2H}\;x\,e^{-x} .
  \label{eq:aeax}
\end{equation}
The function $xe^{-x}$ has its maximum at $x=1$, where it equals $1/e$. Hence
\begin{equation}
  a_{\max} = \frac{V_e^2\,|\sin\gamma_e|}{2\,e\,H},
  \qquad\text{at}\qquad
  \rho_{\text{pk}} = \frac{\beta|\sin\gamma_e|}{H},
  \qquad
  V_{\text{pk}} = \frac{V_e}{\sqrt{e}} .
  \label{eq:allen}
\end{equation}

\paragraph{Independence of ballistic coefficient --- the derivation.}
This is the sharpest statement in \eqref{eq:allen} and it falls straight out of
\eqref{eq:aeax}. The ballistic coefficient $\beta$ enters $a$ \emph{only}
through the dimensionless group $x$; the prefactor $V_e^2|\sin\gamma_e|/(2H)$
contains no $\beta$ at all. And $x$ is a monotone function of density, which
sweeps from $0$ to arbitrarily large values on any entry regardless of $\beta$.
So $x=1$ is always reached, the maximum of $xe^{-x}$ is always attained, and
its value is always $1/e$ times a $\beta$-free prefactor.

What $\beta$ does control is \emph{where} that happens: $x=1$ at
$\rho = \beta|\sin\gamma_e|/H$, so a heavy, slender body (large $\beta$)
penetrates to high density before peaking, and a light, blunt one (small
$\beta$) peaks high in the atmosphere. Physically: every entry body of a given
entry speed and angle experiences the same peak deceleration; what it chooses
by its ballistic coefficient is the altitude at which it takes it, and
therefore how much time it has left afterwards.

\textbf{Measured} at $V_e = 6000~\mathrm{m/s}$, $\gamma_e = -30^{\circ}$,
$H=7200~\mathrm{m}$: closed form $a_{\max} = 459.85~\mathrm{m/s^2}$
($46.89\,g$), $V_{\text{pk}} = 3639.2~\mathrm{m/s}$. The same entry flown at
two ballistic coefficients eight times apart:

\begin{center}
\begin{tabular}{lrr}
\toprule
& $\beta = 8571.4~\mathrm{kg/m^2}$ & $\beta = 1071.4~\mathrm{kg/m^2}$ \\
\midrule
Peak deceleration           & $508.12~\mathrm{m/s^2}$ & $502.16~\mathrm{m/s^2}$ \\
\quad vs.\ closed form      & $+10.50\,\%$            & $+9.20\,\%$ \\
Altitude of the peak        & $4856~\mathrm{m}$       & $19{,}861~\mathrm{m}$ \\
Density at the peak         & $0.624076$              & $0.077647$ \\
\quad vs.\ $\beta|\sin\gamma_e|/H$ & $+4.84\,\%$      & $+4.36\,\%$ \\
Speed at the peak           & $3736~\mathrm{m/s}$     & $3723~\mathrm{m/s}$ \\
$\gamma$ at the peak        & $-31.39^{\circ}$        & $-31.25^{\circ}$ \\
\bottomrule
\end{tabular}
\end{center}

An eightfold change in $\beta$ moves the peak value by $1.2\,\%$ and its
altitude by $15.0~\mathrm{km}$. That is the independence, measured rather than
asserted.

\paragraph{Why the agreement is $10\,\%$ and not $2\,\%$.} Both neglected
effects are quantifiable at the measured peak. Gravity has added speed on the
way down: $3736$ against the predicted $V_e/\sqrt{e} = 3639~\mathrm{m/s}$, i.e.
$+2.7\,\%$, which is $+5.4\,\%$ on $V^2$. And the path has steepened:
$-31.39^{\circ}$ against the assumed $-30^{\circ}$, i.e. $+4.2\,\%$ on
$|\sin\gamma_e|$. Together those predict $+9.8\,\%$; the measured excess is
$+10.50\,\%$. The library's budget is $12\,\%$, and it is \emph{one-sided below
as well}: both neglected effects can only add speed and steepen the path for a
descending ballistic entry, so a simulated peak coming in \emph{under} the
analytic value means too much drag, not a better approximation.

\paragraph{Why the location is compared in density, not altitude.} Altitude is
the logarithm of the predicted quantity. The two cases above disagree by
$4.84\,\%$ and $4.36\,\%$ in density, but the \emph{same} disagreements read as
$6.55\,\%$ and $1.52\,\%$ in altitude --- the constant log offset is being
divided by two different absolute altitudes. A tolerance set in altitude means
different things at different $\beta$; in density it means one thing.

\subsection{Tsiolkovsky}
\label{sec:tsiolkovsky}

In vacuum with gravity off and thrust along the velocity ($\alpha=0$),
\eqref{eq:f4} reduces to $\dot V = T/m$ with $T = T_{\text{vac}}$ constant, and
\eqref{eq:f7} to $\dot m = -\dot m_p$ constant. Then
\begin{equation}
  \frac{\dd V}{\dd m} = \frac{\dot V}{\dot m}
  = -\frac{T}{\dot m_p\,m}
  \;\Longrightarrow\;
  \Delta V = -\frac{T}{\dot m_p}\int_{m_0}^{m_f}\frac{\dd m}{m}
  = \frac{T}{\dot m_p}\ln\frac{m_0}{m_f} ,
\end{equation}
and substituting $\dot m_p = T/(I_{sp}g_0)$ from \eqref{eq:mflow},
\begin{equation}
  \Delta V \;=\; I_{sp}\,g_0\,\ln\frac{m_0}{m_f} .
  \label{eq:tsiolkovsky}
\end{equation}
For the shipped placeholder booster, $m_0 = 32{,}400$~kg,
$m_f = 2{,}400$~kg, mass ratio $13.5$, exhaust speed
$I_{sp}g_0 = 2549.729~\mathrm{m/s}$, giving
$\Delta V = 6636.153368978~\mathrm{m/s}$.

\textbf{Measured:} the propagated vacuum burn gives $6636.153370556$~m/s,
$2.377\times10^{-10}$ relative. The mass history departs from the exact line
$m(t) = m_0 - \dot m_p t$ by $6.8\times10^{-15}$ relative.

\paragraph{This check is blind to every booster constant.} That
$2.4\times10^{-10}$ reads as a very strong result and is a very weak one. Take
the constants in turn:
\begin{center}\small
\begin{tabular}{@{}lp{11.2cm}@{}}
\toprule
Constant & Why it cancels \\
\midrule
$I_{sp}$, $g_0$ & Appear on both sides --- inside
  $\dot m_p = T_{\text{vac}}/(I_{sp}g_0)$ on the simulation side and inside
  $I_{sp}g_0\ln(\cdot)$ on the reference side \\
$T_{\text{vac}}$ & The propagated result is
  $(T_{\text{vac}}/\dot m_p)\ln(m_0/m_f)$; substituting $\dot m_p$ removes it.
  Doubling the thrust halves the burn time and leaves $\Delta V$ untouched \\
Mass ratio & The propagation starts at $m_0$ and terminates on an event set to
  $m_f$; the reference is built from the same fields. A wrong propellant load
  moves both sides identically \\
\bottomrule
\end{tabular}
\end{center}
So the check validates the \emph{shape} of the rocket equation --- that the
integral of $T/m$ really is an exhaust speed times a log mass ratio --- and
nothing whatever about the numbers fed into it. The only defence is to pin each
constant by exact equality at its source, which the library does in
\code{tests/test\_constThrust.m}. A range check is not a pin: bracketing $g_0$
to $|g_0 - 9.80665| < 10^{-3}$ would pass a value wrong in the fourth decimal.

\paragraph{What the ideal figure is worth in practice.} \eqref{eq:tsiolkovsky}
is loss-free. Measured term by term along the shipped ballistic ascent
(89$^\circ$ at the pad pitching to 34$^\circ$ at burnout): gravity loss
$697.23$, drag $48.37$, back pressure $98.85$, steering $33.51$, total
$877.96~\mathrm{m/s}$, leaving a burnout speed of $5768.20~\mathrm{m/s}$ ---
which closes against the propagated value to $0.01~\mathrm{m/s}$. Note that
gravity loss alone is $697~\mathrm{m/s}$; a trajectory steep enough to reach a
1600~km apogee spends most of its burn fighting gravity nearly head-on.

\subsection{Keplerian free-flight range from a burnout state}
\label{sec:kepler}

\paragraph{The development.} Above the sensible atmosphere the flight is
two-body. Take a burnout state $(r_0, V_0, \gamma_0)$ and ask where the
resulting conic meets the impact sphere $r = r_I$.

The specific angular momentum is the radius times the \emph{transverse}
velocity component, which is $V_0\cos\gamma_0$ since $\gamma$ is measured from
the local horizontal:
\begin{equation}
  h = r_0 V_0 \cos\gamma_0 .
  \label{eq:kh}
\end{equation}
The semi-latus rectum and the specific energy follow from $h$ and the vis-viva
form of \eqref{eq:energy}:
\begin{equation}
  p = \frac{h^2}{\mu}, \qquad
  \varepsilon = \frac{V_0^2}{2} - \frac{\mu}{r_0} .
  \label{eq:kpe}
\end{equation}
Eccentricity follows from the energy relation $\varepsilon = -\mu/(2a)$ and
$p = a(1-e^2)$, which give $\varepsilon = -\mu(1-e^2)/(2p)$ and hence
$e^2 = 1 + 2\varepsilon p/\mu = 1 + 2\varepsilon h^2/\mu^2$:
\begin{equation}
  e = \sqrt{1 + \frac{2\varepsilon h^2}{\mu^2}} .
  \label{eq:ke}
\end{equation}
The orbit equation $r = p/(1+e\cos f)$ gives the true anomaly at burnout
through its cosine, and the radial velocity relation
$\dot r = (\mu/h)\,e\sin f$ gives it through its sine:
\begin{equation}
  e\cos f_0 = \frac{p}{r_0} - 1 , \qquad
  e\sin f_0 = \frac{h\,V_0\sin\gamma_0}{\mu} .
\end{equation}
Both are divided by the same positive $e$, and $\operatorname{atan2}$ is
invariant under a common positive scaling, so the eccentricity need not be
divided out:
\begin{equation}
  f_0 = \operatorname{atan2}\!\left(\frac{h V_0\sin\gamma_0}{\mu},\;
        \frac{p}{r_0} - 1\right).
  \label{eq:kf0}
\end{equation}
At the impact radius the orbit equation gives the cosine directly, and the
\emph{descending} branch --- past apogee, on the way down --- is the one with
$f \in (\pi, 2\pi)$:
\begin{equation}
  \cos f_I = \frac{p/r_I - 1}{e} , \qquad
  f_I = 2\pi - \arccos\!\left(\frac{p/r_I - 1}{e}\right).
  \label{eq:kfI}
\end{equation}
The free-flight range angle is the swept true anomaly, and the surface range is
that angle on the impact sphere:
\begin{equation}
  \Delta f = (f_I - f_0)\bmod 2\pi , \qquad
  s_{\text{vac}} = r_I\,\Delta f .
  \label{eq:krange}
\end{equation}

\paragraph{Two guards the implementation needs.} If $|\cos f_I| > 1$ the conic
never reaches the impact radius --- the burnout state is on an orbit whose
perigee is above $r_I$ --- and the formula is meaningless rather than merely
inaccurate; the code raises with the state that produced it. And if $e$ is
numerically zero the true anomaly is undefined; the code refuses below
$10^{-12}$.

\paragraph{Physical reading.} Everything in \eqref{eq:kh}--\eqref{eq:kfI} is
built from $\mu$ and the three burnout numbers, and from nothing about the
vehicle. Free flight does not care what is flying: mass, area and drag
coefficient are all absent. That is exactly why this check is a good one for
the \emph{propagator} and a poor one for the \emph{vehicle}.

\paragraph{Measured, three ways.} The library's own check, inside
\code{run\_ballistic}, re-propagates the burnout state through the full
spherical equations of motion with lift and drag switched off, rotation off and
point-mass gravity \emph{forced} (whatever the user selected), and compares the
propagated great-circle range against \eqref{eq:krange}:
\begin{center}
\textbf{$8.484\times10^{-13}$ relative}, on a $4467.5975~\mathrm{km}$ arc,
against a $10^{-6}$ budget.
\end{center}
Independently, for this note, \eqref{eq:krange} was checked against
\emph{direct two-body integration in Cartesian coordinates} --- a formulation
sharing no line of code and no coordinate system with the spherical
propagator --- at three burnout states, integrating $\ddot{\mathbf{r}} =
-\mu\mathbf{r}/r^3$ with \code{ode89} at $3\times10^{-14}$ relative tolerance
and terminating on a descending crossing of $r = r_I$:

\begin{center}\small
\begin{tabular}{@{}rrrrrr@{}}
\toprule
$r_0$ (km) & $V_0$ (m/s) & $\gamma_0$ (deg) & $e$ & Closed form (km) & Cartesian (km) \\
\midrule
$6480.137$ & $5768.198$ & $47.60$ & $0.800714$ & $4467.597508$ & $4467.597508$ \\
$6458.137$ & $6500.000$ & $25.00$ & $0.510244$ & $7052.123547$ & $7052.123547$ \\
$6498.137$ & $5200.000$ & $40.00$ & $0.772443$ & $3771.430312$ & $3771.430312$ \\
\midrule
\multicolumn{6}{@{}l}{Relative agreement, top to bottom:
  $1.46\times10^{-15}$, \; $1.27\times10^{-14}$, \; $1.91\times10^{-14}$} \\
\bottomrule
\end{tabular}
\end{center}

The first row is the shipped \code{run\_ballistic} burnout state. Note the two
figures are not measuring the same thing: $8.484\times10^{-13}$ is the
\emph{spherical propagator} against the closed form, and includes the
propagator's own integration error; $\le 1.9\times10^{-14}$ is the closed form
against an exact two-body reference, and isolates the algebra of
\eqref{eq:kh}--\eqref{eq:krange}.

\subsection{Great-circle distance and initial bearing}
\label{sec:gc}

\paragraph{Distance.} The haversine form of the central angle between two
points on a sphere:
\begin{equation}
  a = \sin^2\!\frac{\phi_2-\phi_1}{2}
      + \cos\phi_1\cos\phi_2\,\sin^2\!\frac{\lambda_2-\lambda_1}{2},
  \qquad
  d = 2\arcsin\bigl(\min(1,\sqrt{a})\bigr) .
  \label{eq:haversine}
\end{equation}
Multiply by a radius for arc length. The clamp at one is not decoration:
roundoff can push $a$ a hair above unity for near-antipodal points, and
$\arcsin$ would then return a complex number.

Haversine is used in preference to the spherical law of cosines,
\begin{equation*}
  \cos d = \sin\phi_1\sin\phi_2 + \cos\phi_1\cos\phi_2\cos\Delta\lambda ,
\end{equation*}
because for small $d$ that expression evaluates to $1 - d^2/2 + O(d^4)$;
forming it in double precision and then taking $\arccos$ discards about half
the significant digits, leaving an absolute error of order
$\sqrt{\epsilon}\approx 10^{-8}$~rad and returning \emph{exactly zero} below
$d\sim 10^{-8}$~rad. Measured: at a separation of $10^{-10}$~rad, haversine
returns $1.000000\times10^{-10}$ and the cosine form returns exactly $0$.
Terminal-phase ranges are short, so this matters. Haversine's own weak point is the
near-antipodal case, where it degrades to about $\sqrt{\epsilon}$ --- harmless
here.

\paragraph{Initial bearing.} From the spherical triangle formed by the pole and
the two points, the four-parts formula with the cotangents cleared gives
\begin{equation}
  \psi_0 = \operatorname{atan2}\bigl(
    \sin\Delta\lambda\cos\phi_2,\;
    \cos\phi_1\sin\phi_2 - \sin\phi_1\cos\phi_2\cos\Delta\lambda\bigr),
  \qquad \Delta\lambda = \lambda_2-\lambda_1,
  \label{eq:bearing}
\end{equation}
wrapped into $[0,2\pi)$. The first argument is the eastward component of the
initial track direction and the second the northward one, both scaled by the
same positive factor, so their ratio is the tangent of the course. The
\code{atan2} form resolves the full circle without a quadrant test --- which is
where hand-rolled bearing code usually breaks --- and stays well conditioned on
short legs, where numerator and denominator vanish together.

\paragraph{Bearing is not symmetric the way distance is.} Distance satisfies
$d(A,B) = d(B,A)$. Bearing does not satisfy
$\psi_0(B,A) = \psi_0(A,B) \pm \pi$, because the great-circle arc meets the
local meridian at a different angle at each end. Measured, Los Angeles
($33.94^{\circ}$N, $118.41^{\circ}$W) to New York ($40.64^{\circ}$N,
$73.78^{\circ}$W):
\begin{center}
\begin{tabular}{lr}
\toprule
Outbound course, LAX $\to$ JFK      & $65.8669^{\circ}$ \\
Return course, JFK $\to$ LAX        & $273.8410^{\circ}$ \\
Reversing the outbound course would give & $245.8669^{\circ}$ \\
Discrepancy                          & $27.974^{\circ}$ \\
\bottomrule
\end{tabular}
\end{center}
That $28^{\circ}$ is the convergence of the meridians and is the whole point.

\paragraph{But the discrepancy is \emph{not} how you detect a transposition.}
This is worth being precise about, because the natural assumption is the
opposite. Exchanging $\phi_1$ with $\phi_2$ inside \eqref{eq:bearing} moves
\emph{both} courses --- measured, LAX--JFK becomes $86.1590^{\circ}$ and
JFK--LAX becomes $294.1331^{\circ}$ --- and leaves their difference at exactly
the same $27.974^{\circ}$. What a transposition does break is the
\emph{meridional cardinals}, and the reason is structural. On a meridian
$\Delta\lambda = 0$, so the eastward component $\sin\Delta\lambda\cos\phi_2$
of \eqref{eq:bearing} is \emph{exactly} zero --- before and after the swap
alike --- and the transposed answer can therefore only be $0$ or $\pi$,
whichever way the northward component points. Both meridional cardinals invert:

\begin{center}
\begin{tabular}{llrr}
\toprule
Leg & $(\phi_1,\lambda_1)\to(\phi_2,\lambda_2)$ & Correct & $\phi_1\!\leftrightarrow\!\phi_2$ \\
\midrule
Due north & $(0,0)\to(0.1,0)$   & $0.0000^{\circ}$   & $180.0000^{\circ}$ \\
Due south & $(0,0)\to(-0.1,0)$  & $180.0000^{\circ}$ & $0.0000^{\circ}$ \\
Due east  & $(0,0)\to(0,0.1)$   & $90.0000^{\circ}$  & $90.0000^{\circ}$ \\
Due west  & $(0,0)\to(0,-0.1)$  & $270.0000^{\circ}$ & $270.0000^{\circ}$ \\
\bottomrule
\end{tabular}
\end{center}

The north entry is $\operatorname{atan2}(0,\,-\sin 0.1) = \pi$ exactly. The
east--west pair survives untouched because with $\phi_1 = \phi_2 = 0$ there is
nothing to exchange. So the assertion to write is an absolute bearing on a
\textbf{near-meridional} leg, not a reverse-course offset.

The four values are tabulated rather than described because this single claim
has been written down wrongly three times in this project, by three different
authors, the last of them after a review had already established the right
answer. A transposed due-north leg cannot return $90^{\circ}$: that would
require a nonzero eastward component, and on a meridian there is none. The
numbers above are the defence against a fourth attempt.

\paragraph{Three degenerate cases, all conventions rather than answers.}
Equation~\eqref{eq:bearing} hands back a number at three geometries where no
bearing exists. \emph{Coincident} points have no arc: both components vanish,
$\operatorname{atan2}(0,0)$ is defined as zero, and the formula reports due
north for a point to itself. \emph{Antipodal} points admit \emph{every} great
circle, so whatever comes back is an artifact of how two vanishing terms round
--- from $(0,0)$ to $(0,\pi)$ the formula gives exactly $90^{\circ}$, because
$\sin\pi$ evaluates to $1.22\times10^{-16}$ against a true-zero denominator,
and nudging either latitude by $10^{-9}$~rad swings the answer across the whole
circle. And \emph{at a pole} the local north that a bearing is measured from
does not exist: at the north pole every departure direction is south, at the
south pole every departure direction is north, and in both cases longitude ---
and with it any navigational bearing --- is a labelling choice, not a
direction.

Because this quantity is a targeting azimuth, a plausible-looking number at any
of the three is worse than no number, since it flies.
\code{coorbital.util.greatCircleBearing} therefore \emph{raises} rather than
returns at all three: \code{degenerateArc} when
$\operatorname{hypot}$ of the two components --- which is the sine of the
angular separation, so one test catches both the coincident and the antipodal
end --- falls to $10^{-12}$ or below, and \code{polarOrigin} when
$|\cos\phi_1|\le 10^{-12}$, which needs its own test because the components do
not vanish at a pole. Near-antipodal geometry stays ill-conditioned well
outside the refusal band; the guard removes only the cases that have no answer
at all.

\subsection{Summary of measured agreements}

\begin{center}\small
\begin{tabular}{@{}llr@{}}
\toprule
Check & Budget & Measured \\
\midrule
Vacuum specific energy, 300~s lofted arc          & $<10^{-8}$ rel & $1.281\times10^{-13}$ \\
Rotating-frame Jacobi integral, same arc          & $<10^{-8}$ rel & $1.107\times10^{-13}$ \\
\quad \emph{same, with the Coriolis sign flipped} & --- & $1.468\times10^{-13}$ \\
Equilibrium glide $V(r)$, 45 $\to$ 20~km          & $<3\,\%$ & $1.51\,\%$ \\
Allen--Eggers peak, $\beta$ 8$\times$ apart       & $0<\text{rel}<12\,\%$ & $+10.50\,\%$, $+9.20\,\%$ \\
Allen--Eggers peak location, in density           & $<6\,\%$ & $+4.84\,\%$, $+4.36\,\%$ \\
Peak value spread over 8$\times$ in $\beta$       & $<3\,\%$ & $1.2\,\%$ \\
\code{boost3DOF} $\to$ \code{glide3DOF}, engine dead & $<10^{-12}$ & $0.000\times10^{0}$ \\
Thrust force increment, $\alpha,\sigma\ne 0$      & $<10^{-12}$ & $3.642\times10^{-16}$ \\
Tsiolkovsky, full vacuum burn                     & $<10^{-8}$ rel & $2.377\times10^{-10}$ \\
Mass linearity over the burn                      & $<10^{-9}$ rel & $6.821\times10^{-15}$ \\
Keplerian range, propagated vs closed form        & $<10^{-6}$ rel & $8.484\times10^{-13}$ \\
Keplerian range, closed form vs Cartesian two-body & --- & $\le1.91\times10^{-14}$ \\
Gravity/aero projections vs vector reconstruction & --- & $\le9.5\times10^{-16}$ \\
Thrust projections vs vector reconstruction       & --- & $\le1.6\times10^{-16}$ \\
Rotating-Earth terms vs vector reconstruction     & --- & $\le1.1\times10^{-13}$ \\
\bottomrule
\end{tabular}
\end{center}

\paragraph{A caution on the bit-exact row.} The
\code{boost3DOF}$\,\to\,$\code{glide3DOF} reduction measures $0.000\times10^0$
because every shared term in the powered file is a character-identical copy of
the unpowered expression evaluated in the same order, and with $T=0$ the
perturbations $(0 - a_D)$ and $(0 + a_L)$ are IEEE-exact identities. The check
therefore \emph{proves} that the six shared equations are the glide equations
--- no transcription drift, no dropped term, no altered sign is possible --- and
is simultaneously \emph{zero-bit evidence of correctness}, because any error in
the glide equations is inherited and cancels invisibly. It validates identity,
not physics. A reduction test validates only the terms that survive the
reduction.

%=====================================================================
\section{Targeting}
\label{sec:targeting}

Given a launch point and a destination rather than a launch point and an
azimuth, two unknowns must be found. Over a non-rotating sphere they decouple
completely, and the library solves them as two independent problems.

\subsection{The azimuth is closed-form}

The initial bearing of the launch-to-target great circle, \eqref{eq:bearing},
\emph{is} the heading state $\psi$ --- both are measured clockwise from north
--- so it goes straight into $\mathbf{x}_6$ with no conversion, no iteration and
no tolerance. For the shipped geometry (Hawaii to the California coast) it is
$\psi_0 = 0.988331159~\mathrm{rad} = 56.627204^{\circ}$, and the flown initial
heading is asserted equal to it to $10^{-13}$.

\paragraph{Exact only because the Earth does not rotate here.} With
$\omega = 0$ the ground track of a zero-bank trajectory is a great circle, so
the initial bearing of the launch-to-target arc is the whole answer. Switch the
Earth rate on and the target is carried east under the vehicle for the entire
flight --- of order 600~km at mid latitude over a 27-minute flight --- so the
vehicle must be aimed not at where the target is but at where it \emph{will
be}. The azimuth then depends on the flight time, the flight time on the
cutoff, and the cutoff on the azimuth: what was a closed form becomes an outer
iteration wrapped around the range solve. The library does not have that outer
loop, and refuses to pretend it does: with rotation enabled it prints a caution
and reports the miss against a target that stood still.

\subsection{The range solve is a bisection on thrust-termination time}

The boost phase already ends at $t_{\text{span}}(2)$ when the burnout event has
not fired first, so a shortened span \emph{is} a thrust cutoff and no new
machinery is needed. Less burn gives less energy gives less range: the map
$t_{\text{cut}} \mapsto R$ is monotonic and single-valued, which is exactly
what makes bisection safe.

Bisection, rather than a secant or Newton method, because each evaluation is a
full trajectory propagation of several seconds, the derivative is unavailable,
and the function is only piecewise smooth. Bisection never leaves its bracket,
needs no derivative, and cannot be thrown off by a local kink; it costs a fixed
number of evaluations per decade of tolerance, which is the right trade when
the alternative occasionally has to be restarted. Convergence is tested on the
\emph{achieved value} --- a kilometre of range --- not on the parameter, because
that is the quantity the user cares about.

\paragraph{Why cutoff time and not something else.} The two obvious
alternatives both bifurcate.
\begin{itemize}
\item \textbf{Loft angle} has a lofted and a depressed branch either side of a
maximum-range hump --- two cutoffs per achievable range --- so a root finder
lands on whichever branch it started nearest. Monotonicity is what bisection
assumes and does not check; with several crossings it finds one and does not
say which.
\item \textbf{Glide handoff altitude} bifurcates on the skip phugoid's troughs.
Measured in this library: moving the handoff from 15~km to 30~km costs
$1881~\mathrm{km}$ of range --- a quarter of the flight --- \emph{with every
phase still reporting nominal}, because every event fired exactly as asked.
\end{itemize}
Thrust termination is also what real systems do for energy management. Its
price is that an early cutoff throws away unburned propellant with the
jettisoned booster, so the summary reports how much went overboard; at a deep
cutoff that is most of the load.

\textbf{Measured} on the shipped configuration: required $3811.2399$~km,
achieved $3811.7512$~km, residual $511.24$~m, in 10 bisection iterations, at a
cutoff of $75.681975~\mathrm{s}$ ($93.99\,\%$ of the full burn). The reachable
band, which is the two bracket evaluations the solve makes anyway, is
$202.861$ to $7737.630$~km. A target outside that band is not an error: the
solver returns \code{converged = false} with the band, and the caller phrases
the refusal in its own vocabulary. What must never happen is a silent nearest
miss, because a nearest miss \emph{looks} like a solution.

\subsection{Cross-range is a property of the configuration, not a guarantee}

The bisection matches the great-circle \emph{distance} from launch to impact.
Nothing in it steers the track sideways. Cross-range comes out at zero in the
shipped configuration because the bank angle is zero throughout, and a
zero-bank trajectory over a non-rotating sphere stays on the great circle it
left on --- which is a statement about $\sigma=0$ and $\omega=0$ together, not
a general property of the solver.

The library therefore \emph{measures} it rather than asserting it. The
cross-track offset of the impact point from the launch-to-target great circle
is computed by the standard spherical relation
\begin{equation}
  x_{\text{track}} \;=\; r_I \arcsin\!\bigl(
     \sin\Delta\theta_{\text{fly}}\;
     \sin(\psi_{\text{fly}} - \psi_{\text{launch}})\bigr),
  \label{eq:xtrack}
\end{equation}
in which \emph{all three inputs are quantities at the launch point}:

\begin{itemize}
\item $\Delta\theta_{\text{fly}}$ is the great-circle central angle from the
launch point to the \emph{achieved impact} point, \eqref{eq:haversine} on the
flown terminal position;
\item $\psi_{\text{fly}} = \operatorname{bearing}(\text{launch},
\text{impact})$ is the \emph{initial} bearing of the launch-to-impact arc,
\eqref{eq:bearing} evaluated with the flown impact point as its second
argument;
\item $\psi_{\text{launch}} = \operatorname{bearing}(\text{launch},
\text{target}) = \psi_0$ is the initial bearing of the intended launch-to-target
arc --- the same closed form that seeded $\mathbf{x}_6$ in
Section~\ref{sec:targeting}.
\end{itemize}

\noindent
The difference of the two bearings is wrapped into $(-\pi,\pi]$ before the
sine, which fixes the sign convention as positive to the right of the intended
course. In the code these are \code{angFly}, \code{psiFly} and
\code{psiLaunch}, and \code{psiFly} is a second call to
\code{greatCircleBearing} from the launch point --- not a state read out of the
trajectory.

\paragraph{Why it is not the terminal heading, which is the plausible mistake.}
$\psi_{\text{fly}}$ is emphatically \emph{not} $\mathbf{x}_6$ at impact. The
terminal heading is tangent to the flown track at its \emph{far} end, and it
differs from the initial launch-to-impact bearing by exactly the meridian
convergence of Section~\ref{sec:gc} --- the same effect that puts
$27.974^{\circ}$ between the LAX--JFK outbound and return courses. Substituting
it into \eqref{eq:xtrack} leaves a formula that still looks right, still
returns a number of plausible magnitude, and is wrong by tens of kilometres on
an intercontinental arc. Both angles in \eqref{eq:xtrack} have to be anchored
at the launch point, because the relation is a statement about the spherical
right triangle standing on the intended great circle \emph{there}: one leg
along the intended course, the hypotenuse along the launch-to-impact arc, and
the included angle between them.

The offset is printed every run, with a warning when it exceeds the range
tolerance. Measured:

\begin{center}
\begin{tabular}{lrr}
\toprule
& Shipped, $\sigma = 0$ & $\sigma = 75^{\circ}$ descent \\
\midrule
Cross-track offset & $-4.06\times10^{-7}$~m & $21{,}515.22$~m \\
Total miss         & $511.24$~m & $21{,}524.70$~m \\
Range residual     & $511.24$~m & $-585.20$~m \\
\bottomrule
\end{tabular}
\end{center}

At $75^{\circ}$ of bank the range residual is still under a kilometre --- the
bisection did its job --- and the vehicle lands 21.5~km to the side.
Closing that requires an outer iteration on azimuth, which the library does not
have. The honest reading is that the shipped script solves \emph{downrange
only}, and the cross-track line in its summary is the evidence for the
zero-cross-range claim rather than a restatement of it.

%=====================================================================
\section{Limitations}

Stated plainly, because the model's boundaries are part of its specification.

\begin{itemize}
\item \textbf{Spherical Earth only.} No $J_2$, no ellipsoid, no geodetic
altitude. The interfaces are shaped so that an oblate gravity model drops in
without editing the equations of motion, and the $g_\phi$ terms exist for
exactly that day, but they are unexercised today.
\item \textbf{The rotating-Earth terms are off by default} and are exercised by
the test suite but by no shipped trajectory. The targeting azimuth is
closed-form only with them off.
\item \textbf{The $g_\phi$ terms are verified only by this note}, using an
injected synthetic gravity model. Nothing in the committed test suite can see
them.
\item \textbf{Aerodynamics are constant.} $\alpha$ does not reach $C_L$, so
during unpowered flight the only live control is bank angle. Mach number is
computed and ignored.
\item \textbf{The atmosphere is isothermal-exponential} and, as
Section~\ref{sec:models} quantifies, not exactly hydrostatic under
inverse-square gravity. Its purpose is closed-form validation, not fidelity.
\item \textbf{No aerothermodynamics.} \code{noseRadius} is carried on the
vehicle struct as a placeholder for a deferred Sutton--Graves model; nothing
reads it.
\item \textbf{No 6-DOF, no actuator dynamics, no autopilot.} Commanded angles
are achieved instantaneously.
\item \textbf{Every vehicle and booster number is a placeholder.} The booster's
structural coefficient in particular, $1500/31500 = 4.76\,\%$, is optimistic by
roughly a factor of two against real large solids; a realistic $11.2\,\%$ stage
with the same propellant load would deliver about $5.1$ rather than
$6.6~\mathrm{km/s}$ of ideal $\Delta V$. The alignment between this booster's
burnout speed and the 60~km, 6~km/s entry interface used elsewhere in the
library is an artefact of that optimism, not a design expectation.
\end{itemize}

%=====================================================================
\begin{thebibliography}{9}

\bibitem{Vinh1980}
Vinh, N.X., Busemann, A., Culp, R.D.,
\emph{Hypersonic and Planetary Entry Flight Mechanics},
University of Michigan Press, 1980. Eqs.~(2.28)--(2.33) are the flight-mechanics
equations of Section~\ref{sec:dynamics}; Sec.~5.3 the equilibrium glide;
Ch.~6 the Allen--Eggers entry.

\bibitem{Allen1958}
Allen, H.J., Eggers, A.J.,
``A Study of the Motion and Aerodynamic Heating of Ballistic Missiles Entering
the Earth's Atmosphere at High Supersonic Speeds,''
NACA TR-1381, 1958.

\bibitem{Eggers1958}
Eggers, A.J., Allen, H.J., Neice, S.E.,
``A Comparative Analysis of the Performance of Long-Range Hypervelocity
Vehicles,'' NACA TR-1382, 1958.

\bibitem{Sanger1944}
S\"anger, E., Bredt, I.,
``A Rocket Drive for Long Range Bombers,''
Deutsche Luftfahrtforschung UM-3538, 1944.

\bibitem{Bate1971}
Bate, R.R., Mueller, D.D., White, J.E.,
\emph{Fundamentals of Astrodynamics}, Dover, 1971. Ch.~6, the free-flight range
angle of a ballistic missile.

\bibitem{Sutton2017}
Sutton, G.P., Biblarz, O.,
\emph{Rocket Propulsion Elements}, 9th ed., Wiley, 2017.
Sec.~3.2--3.3 (choked flow and the thrust pressure term), Sec.~4.1 (the ideal
rocket equation), Ch.~12 (representative solid-motor performance).

\bibitem{Sinnott1984}
Sinnott, R.W., ``Virtues of the Haversine,''
\emph{Sky and Telescope}, Vol.~68, No.~2, 1984, p.~159.

\bibitem{Bowditch}
Bowditch, N., \emph{The American Practical Navigator},
Pub.~No.~9, National Geospatial-Intelligence Agency, chapter on Great-Circle
Sailing.

\bibitem{NumericalRecipes}
Press, W.H., Teukolsky, S.A., Vetterling, W.T., Flannery, B.P.,
\emph{Numerical Recipes}, 3rd ed., Cambridge University Press, 2007, Sec.~9.1
(bisection).

\bibitem{WGS84}
\emph{Department of Defense World Geodetic System 1984},
NIMA TR8350.2, 3rd ed., 1997. Source of $\mu$, $r_E$ and $\omega$.

\end{thebibliography}

\end{document}
